\input hofstadter.opm

Notes on the Cover 
A Spontaneous Essay on Whirly Art and Creativity 
The drawing on the cover is a somewhat atypical example of a non- representational form of art I devised and developed over a period of years quite a, long time ago, and which my sister Laura once rather light-heartedly dubbed "Whirly Art". The name stuck, for better or for worse. Generally speaking,. I did Whirly Art on long thin strips of paper (available in rolls for adding machines) rather than on sheets of standard format. A typical piece of Whirly Art is five or six inches high and five or six feet long. Many are ten feet long, however, and some are as much as fifteen or even twenty feet in length. The one-dimensionality of Whirly Art was deliberate, of course: I was inspired by music and drew many visual fugues and canons. The time dimension was replaced by the long space dimension. I used the narrow width of the paper to represent something like pitch (although there was no strict mapping in any sense). A "voice" would be a single line tracing out some complex shape as it progressed in "time" along the paper. Several such voices could interact, and notions of what made "good" or "bad" visual harmony' or counterpoint soon became intuitive to me. The curvilinear motions constituting a single voice came from a blend of alphabets. At that time (the mid-60's), I was absolutely fascinated by the many- writing systems found in and around India, exemplified by Tamil Sinhalese, Kanarese, Telugu, Bengali, Hindi, Burmese, Thai, and many others. I studied some of them quite carefully, and even invented one of my own, based on the principles that most Indian scripts follow. It was natural that the motions my hand and mind were getting accustomed to would find their way into my visual fuguing. Thus was born Whirly Art. Over the next, several years, I did literally thousands of pieces of Whirly Art. Each one was totally improvised-in pen-so that there was no going back, a mistake was a mistake! Alternatively, a mistake could be interpreted as a very daring move from which it would be difficult, but not impossible to recover gracefully. In other words, what seemed at first to be a disastrous mistake could turn into a joyful challenge! (I am sure that jazz improvisers will know exactly what I am talking about.) Sometimes, of course, I would 



Notes on the Cover                                                               XVIII Fail, but other times I would succeed (at least by my own standards, since I was both performer and "listener"). Whirly Art became a (very) highly idiosyncratic language, with its own esthetic and traditions. However, traditions are made to be broken, and as soon as I spotted a tradition, I began experimenting around, violating it in various ways to see how I might move beyond my current state-how I might "jump out of the system". Style succeeded style, and I found myself paralleling the development of music. I moved from baroque Whirly Art (fugues, canons, and so forth) to "classical" Whirly Art, thence to "romantic" Whirly Art. After several years (it was now the late 60's), I reached the twentieth century, and found myself spiritually imitating such favorite composers of mine as Prokofiev and Poulenc. I did not copy any pieces specifically, but simply felt a kinship to those composers' style. Whirly Art iS not translated music, but metaphorical music. It is natural to wonder if I managed to jump beyond the twentieth century and make visual 21st-century music. That would have been quite a feat! Actually, in the early 70's I found that I simply was slowing down in production of Whirly .Art. It had taken me seven years to recapitulate the history of Western music! At that point, I seemed to run out of creative juices. Of course, I could still make new Whirly Art then, as I can now-but I simply was less often inclined to do so. And today, I hardly ever do any Whirly Art, although the way that I draw curvy lines and letterforms bears indelible marks of Whirly Art. The piece on the cover, then, is atypical because it was done on an ordinary sheet of paper and has no direction of temporal flow. Also, the really is no concept of counterpoint-in it. Still, it has something of a Whirly Art spirit. There are also seven Whirly alphabets, in the book, one on each of the title pages of the seven sections. They ;are all somewhat atypical as, well, but for slightly different reasons. Each was done on an ordinary sheet of paper but there is still always a clear flow, namely from `A' to, `Z': The real atypicality is the fact that genuine letters from a genuine alphabet are being used. I usually eschewed real letters, preferring-to use shapes inspired by letters-shapes more complex and, well, "whiny" than most letters, even more so than Tamil or Sinhalese letters, which are pretty, darn whiny. Whirly Art is, I feel, quite possibly the most creative thing I have ever done. That, of course, is my opinion. Other people may disagree. It is a fairly strange and idiosyncratic form of art, however, and cannot be instantly understood. It has its own logic, related to the logics of musical harmony and counterpoint, Indian alphabets, gestalt perception, and who knows what else. I've kept it all quite literally in my closet for years-rolled up and piled into many paper bags and cardboard boxes. Because of its physical awkwardness, it is hard to show to people. But Whirly Art itself, and the experience of doing it, is an absolutely central fact about my way of looking at art, music, and creativity. Practically every time I write about creativity, 



Notes on the Cover                                                                   XIX some part of my mind is re-enacting Whirly Art experiences In other words a lot of my convictions about creativity come from self-observation rather than from scholarly study of the manuscripts or sketches of various composers or painters or writers or scientists. Of course, I have done some of that type of scholarship too, because I am fascinated by creativity in general-but I feel that to some extent "you don't really understand it unless you've done it", and so I rely a great deal on that personal experience. I feel that way that "I know what I'm talking about." However, I would make a slightly stronger statement: Any two creative things that I've done seem to be, at some deep level, isomorphic. It's as if Whirly Art and mathematical discoveries and strange dialogues and little pieces of piano music and so on are all coming from a very similar core, and the same mechanisms are being exploited over and over again, only dressed up differently. Of course it's not all of the same quality: my real music-is not as good as my visual music, for instance. But because I have this conviction that the core creativity behind all these things is really the same (at least in my own case), I am trying like mad to get at, and to lay bare, that core. For. that reason I pursue ever-simpler domains in which I can feel myself doing "the same thing". In Chapter 24 of this book-in some sense the most creative Chapter, not surprisingly-I write about three of those domains, the Seek-Whence domain, the Copycat domain, and the Letter Spirit domain. It is the Letter Spirit domain--"gridfonts" in particular-that is currently my most intense obsession. That domain came out of a lifelong fascination with our alphabet and other writing systems. I simply boiled away what I considered to be less interesting aspects of letterforms-I boiled and boiled until I was left with what might be called the "conceptual skeletons" of letterforms. That is what gridfonts are about. People who have not shared my alphabetic fascination often underestimate at first the potential range of gridfonts, thinking that there might be a few and that's all. That is dead wrong Thee are a huge number of them, and their variety is astounding. As I look at the gridfonts I produce-and as I feel myself producing a gridfont I feel that what I am doing is just Whirly Art all over again, in a new and ridiculously constrained way. The same mechanisms of 'shape transformation, the same quest for grace and harmony, the same intuitions bout what works and what doesn't, the same desire to "jump out of the system"-all this is truly the same. Doing gridfonts is therefore very exciting me and provides a new proving ground for my speculations. The one advantage that gridfonts have over Whirly Art is that they are preposterously constrained. This means that the possibilities for choice can be watched much more easily. It does not mean that a choice can be explained easily, but at least it can be watched. In a way, gridfonts are allowing me re-experience the Whirly-Art period of my life, but with the advantage several years' thinking about artificial intelligence and how I would like t try to make it come about. In other words, I can now hope that perhaps I 



Notes on the Cover                                                                   XX Can get a Handle-a bit of one, anyway-on w at is going on in creativity by means of computer modeling of it. Since I feel that in a fundamental sense, Whirly-Art creativity is no deeper, than gridfont creativity, the study of gridfont creation-more specifically, the computer modeling of gridfont creation-could reveal some things that ' I have sought for a long time. Therefore the next few years will be an important time for me-a time to see if I can really get at the essence, via modeling, of what my mind is doing when I create something that to me is , excitingly novel. This book, as it says on its cover and in the Introduction, deals with Mind and Pattern. To me, boiling things down to their conceptual skeletons is the royal road to truth (to mix metaphors rather horribly). I think that a lot of truth about Mind and Pattern lies waiting to be extracted in the tiny domains that I have carved out very painstakingly over the past seven years or so in Indiana. I urge you to keep these kinds of things in mind as you read this book. This "confession", coming as it does in a most unexpected place, is a very spontaneous one and probably captures as well as anything could the reason that my research is focused as it is, and the reason that I wrote this book. 



Notes on the Cover                                                                 XXI Introduction 
This book takes its title from the column I wrote in Scientific American between January 1981 and July 1983. In that two-and-a-half-year span, I produced 25 columns on quite a variety of topics. My choice of title deliberately left the focus of the column somewhat hazy, which was fine with me as well, as with Scientific American. When Dennis Flanagan, the magazine's editor, wrote to me in mid-1980 to offer me the chance to write a column in that distinguished publication, he made it clear that what was desired was a bridge between the scientific and the literary viewpoints, something he pointed out Martin Gardner had always done, despite the ostensibly limiting title of his column, "Mathematical Games". Here is how Dennis put it in his letter: - 
I might emphasize the flexible nature of the department we have been calling "Mathematical Games". As you know, under this, title, Martin has written a great deal that is neither mathematical nor game-like. Basically, "Mathematical Games", has been Martin's' column to talk about-anything under the sun that interests him. Indeed, in our view, the main import of the column has been to demonstrate that a modern intellectual can have a range of interests that are confined by such words as "scientific" or "literary". We hope that whoever succeeds Martin will feel free to cover his own broad range of interests, which re unlikely to be identical to Martin's. 
What a refreshingly open attitude! So I was being asked to be the successor to' Martin Gardner-but not necessarily to continue the same column, Rather than filling the same role as Martin had, I would merely occupy the same physical spot in the magazine. I had been offered a unique opportunity to say pretty much anything I wanted to say to a vast, ready-made audience, in a prestigious context. Carte blanche, in short. What more could I ask? Even so, I had to deliberate long' and hard about whether to take it, because I did not consider myself primarily a writer, but a thinker and researcher, and time taken in writing would surely be time taken away from research. The conservative pathway, following what was known, would have been to say no, and just do research, The adventurous pathway, exploring the new opportunity and forsaking 



Introduction                                                                       xxii some research, was tempting. Both were risky, since I knew that, either way I would inevitably wonder, "How would things have gone had I decided the other way?" Moreover, I had no idea how long I might write my column, since that was not stipulated. It, could go on for many years-or I could, decide it was too much for me, and quit after a year. In a way, I knew from the beginning that I would take the offer, I guess because I am basically more adventurous than I am conservative. But it was a little like purchasing new clothes: no matter how much you like them, you still want to see how you look in them before you buy them, so you put them on and parade around the store, looking at yourself in the mirror and asking whoever is with you what they think of it. So I talked it over with numerous people, and finally decided as I had expected: to take the offer. 
\separator

For the first year, Martin Gardner and I alternated columns. I have to, admit that even though I was utterly free to "be myself", I felt somewhat. tradition-bound. True, I had metamorphosed his title into my own title (see Chapter 1 for an explanation), but I was aware that readers of Martin's column would, naturally enough, be expecting a similar type'1 of fare. It took a little while for me to test the waters, getting reader reactions and seeing if the magazine was satisfied with my performance, a performance very different in style from Martin's, after all. Needless to say, some readers were "disappointed that I was not a clone of Martin Gardner, but others complimented me on how I had managed to keep the same level of quality while changing the style and content greatly. It was hard, knowing. that people were constantly comparing me with someone very different from me. It was particularly. hard when people who should have known better really confused my role with Martin's. For instance; as late as June 1983, at a conference on artificial intelligence, a colleague who-spotted me came up to me and eagerly told me a math puzzle he'd just discovered and solved, hoping I would put it in my "Mathematical Games" column,. How often did I have to tell people that my column was not called "Mathematical Games" I doubt that anyone loved Martin Gardner's column more than I did, or owed more to it. Yet I did not want my identity confused with someone else's. So writing this column and being in the shadow of, someone superlative was not always easy. But I think I hit my stride and comfortable with, my new role after a few months, In 1982, -Martin retired, leaving the space entirely to me. It was -a chore to be sure, to get a column out each month, but it was also a lot of fun. In any case, what mattered to me the most was to do my best to make the column interesting and diverse and highly provocative. I took Dennis' offer quite literally, not restricting myself to purely scientific topics, but venturing into musical and literary topics as well. After a year and a half, I was beginning to wonder how long I could sustain 



Introduction                                                                          xxiii it without seriously jeopardizing my research. I decided to divide up my long list of prospective topics into categories: columns I would love to do, columns I would simply enjoy doing, and columns I could write with interest but no real passion. I found I had about a year's worth left in the first category, maybe another year's worth left in the second, and then a large number in the third. It seemed, then, that in another year or so it would be a good time to reassess the whole issue of writing the column. As it turned out, my thinking was quite consonant with evolving desires at the editorial level of the magazine. They were most interested in launching a new column to be devoted to the recreational aspects of computing, and our plans dovetailed well. My column could be phased out just as the new one was being phased in. And that is the way it came to pass, with two surprise columns by Martin Gardner filling the gap. My farewell to readers came as a postscript to Martin's final column, in September 1983. Thus my era as a columnist came to an end. As I look back on it, I feel it lasted just about the right length of time: long enough to let me get a significant amount said, but not so long that it became a real drag on me. This way, at least, I got to explore that avenue that was so tempting, and yet it didn't radically alter the course of my life. So in sum, I am quite pleased with my stint at Scientific American. I am proud to have been associated with that venerable institution, and to have filled that unique slot for a ,time, especially coming right on the heels of someone of such high caliber. 
\separator

The diversity of my columns is worth discussing for a moment. On the surface, they seem to wander all over the intellectual map=-from sexism to music to art to nonsense, from game theory to artificial intelligence to molecular biology" to the Cube, and more. But there is, I believe, a deep underlying, unity to my columns. I felt that gradually, as I wrote more and o of-them, regular readers would start to see the links between disparate ones, so that after a while, the coherence of the web would be quite clear. My image of this' was always geometric. I envisioned my intellectual "home territory" as a rather large region in some conceptual space, a region that most people do not see as a connected unit. Each new column was in a way a new "random dot" in that conceptual space, and as dots began peppering the space more fully over the months, the shape of my territory would begin to emerge more clearly. Eventually, I hoped, there would emerge' a clear `region associated with the name "Metamagical Themas". Of course I wonder if my 25 1/2 columns are sufficient to convey the connectedness of my little patch of intellectual territory, or if, on the .contrary, they would leave a question mark in the mind of someone who read them all in succession without any other explanation. Would it simply seem like a patchwork quilt, a curious potpourri? Truth to tell, I suspect that 5 columns are not quite enough, on their own. Probably the dots are' too 



Introduction                                                                         xxiv sparsely distributed to suggest the rich web of potential cross-connections there. For that reason, in drawing all my columns together to form a book; decided to try to flesh out that space by including a few other recent 4rritings of mine that might help to fill some of the more blatant gaps. There are seven such pieces included (indicated by asterisks in the table of contents). I believe they help to unify this book. If someone were to ask me, "What is your new book about, in a word?", 1 "Would probably mutter something like "Mind and Pattern". That, in fact, was one title I considered for the column, way back when. Certainly it tells what most intrigues me, but it doesn't convey it quite vividly or passionately enough. Yes, I am a relentless quester after the chief patterns of the universe--central organizing principles, clean and powerful ways to categorize what is out there". Because of this, I have always been pulled to mathematics. Indeed, even though I dropped the idea of being a professional mathematician many years ago, whenever I go into a new bookstore, I always e a beeline for the math section (if there is one). The reason is that I feel that mathematics, more than any other discipline, studies the fundamental, pervasive patterns of the universe. However, as I have gotten older, I have come to see that there are inner mental patterns underlying our ability to conceive of mathematical ideas, universal patterns in human minds that make them receptive not only to the patterns of mathematics but *'also to abstract regularities of all sorts in the world. Gradually, over the years my focus of interest has shifted to those more subliminal patterns of memory and associations, and away from the more formal, mathematical ones. Thus my interest has turned ever more to Mind, the principal apprehender of pattern, as well as the principal producer of certain kinds of pattern. . To me, the deepest and most mysterious of all patterns is music, a product of the mind that the mind has not come close to fathoming yet. In some sense all my research is aimed at finding patterns that will help us to understand the mysteries of musical and visual beauty. I could be bolder and say, `''I seek to discover what musical and visual beauty really are." However, I don't believe that those mysteries will ever be truly cleared up, nor do I wish them to be. I would like to understand things better, but I don't want to understand them perfectly. I don't wish the fruits of my research to include a mathematical formula for Bach's or Chopin's music. Not that I think it possible. In fact, I think the very idea is nonsense. But even though I find the prospect repugnant, I am greatly attracted by the effort to do as do as much as possible in that direction. Indeed, how could anyone hope to approach the concept of beauty without deeply studying the nature of formal patterns and their organizations and relationships to Mind? How can anyone fascinated by beauty fail to be intrigued by the notion of a "magical formula" behind it all, chimerical though the idea certainly is? And in this day and age, how can anyone fascinated by creativity and beauty fail to "see in computers the ultimate tool for exploring their essence? Such ideas are 



Introduction                                                                           xxv thee inner fire that propels my research and my writings, and they are the core of this book. There is another aspect of my inner fire that is brought out in the writings here collected, particularly toward the end, but it pops up throughout. That is a concern with the global fate of humanity and the role of the individual in helping determine it. I have long been an activist, someone who periodically gets fired up by some cause and ardently works for it, exhorting everyone else I come across to get involved as well. I am a fierce believer in the value of passion and commitment to social causes, someone baffled and troubled by apathy. One of my personal mottos is: "Apathy on the individual level translates into insanity at the mass level", a saying nowhere better exemplified than by today's insane dedication of so many human and natural resources to the building up of unimaginably catastrophic arsenals, all while mountains of humanity are starving and suffering in horrible ways. Everyone knows this, and yet the situation remains this way, getting worse day by day. We do live in a ridiculous world, and I would not wish to talk about the world without indicating my confusion and sadness, but also my vision and hope, concerning our shared human condition. 
\separator

Inevitably, people will compare this book with my earlier books, Gödel Escher, Bach: an Eternal Golden Braid and The Mind's I, coedited with my friend Daniel Dennett. Let me try for a moment to anticipate them. GEB was a; unique sort of book-the detailed working-out of a single potent spark. It was a kind of explosion in my mind triggered by my love with mathematical logic after a long absence. It was the first had tried to write anything long, and I pulled out all the stops. In particular, I made a number of experiments with style, especially in writing dialogues .based on musical forms such as fugues and canons. In essence,'' GEB was one extended flash having to do with Kurt Gödel’s famous incompleteness theorem, the human brain, and the mystery -'consciousness. It is well described on its cover as "a metaphorical fugue on minds and machines". The Mind's 1 is very different from Gödel, Escher, Bach. It is an extensively annotated anthology rather than the work of a single person. It is far more like a monograph than GEB is, in that it has a unique goal: to probe the mysteries of 'matter and consciousness in as vivid and jolting a -way as possible, through stories that anyone can read and understand, followed by careful commentaries by Dan Dennett and myself. Its subtitle is "Fantasies and Reflections on Self and Soul". One thing that GEB and The Mind's I have in common is their internal structure of alternation. GEB alternates between dialogues and chapters; while The Mind's I alternates between fantasies and reflections. I guess I like 



Introduction                                                                        xxvi This contrapuntal mode, because it crops up again in the present volume. Here, I alternate between articles and postscripts. If GEB is an elaborate fugue on one very complex theme, and MI is a collection of-many variations on a theme, then perhaps MT is a fantasia employing several themes. If it were not for the postscripts, I would say that it was disjointed. However, I have made a great effort to tie together the diverse themes-Themas-by writing extensive commentaries that cast the ideas of each article in the light of other articles in the book. Sometimes the postscripts approach the length of the piece they are "post", and in one case (Chapter 24) the postscript is quite a bit longer than its source. 'The reason for that particularly long postscript is that I decided to use it to describe some aspects of my own current research in artificial intelligence. There are other places as well in the book where I touch on my research ideas, though I never go into technical details. My main concern is to give a clear idea of certain central riddles about how minds work, diddles that I have run across over and over again in different guises. The questions I raise are difficult but I find them as beguiling as mathematical ones. In any case, this book will give readers a better understanding of how. my research and the rest of my ideas fit together. 
\separator

One aspect of this book that, I must admit, sometimes makes me uneasy the striking disparity in the seriousness of its different topics. How can both Rubik's Cube and nuclear Armageddon be discussed at equal length one book by one author? Partly the answer is that life itself is a mixture things of many sorts, little and big, light and serious, frivolous and formidable, and Metamagical Themas reflects that complexity. Life is not worth living if one can never afford to be delighted or have fun. There is another way of explaining this huge gulf. Elegant mathematical structures can be as central to a serious modern worldview as are social concerns, and can deeply influence one's ways of thinking about anything ---even such somber and colossal things as total nuclear obliteration. In er to comprehend that which is incomprehensible because it is too huge too complex, one needs simpler models. Often, mathematics can provide right starting point, which is why beautiful mathematical concepts o pervasive in explanations of the phenomena of nature on the microlevel. They-are now proving to be of great help also on a larger scale, as Robert Axelrod's lovely work-on the Prisoner's Dilemma so impeccably demonstrates (see \chapter {29}
The Prisoner's Dilemma is poised about halfway between the Cube and Armageddon, in terms of complexity, abstraction, size, and seriousness. I submit that abstractions of this sort are direly needed in our times, because many people-even remarkably smart people-turn off when faced with issues that are too big. We need to make such issues graspable. To make 



Introduction                                                                         xxvii them graspable and fascinating as well, we need to entice people with the beauties of clarity, simplicity, precision, elegance, balance, symmetry, and so on. Those artistic qualities, so central to good science as well as to good insights about life, are the things that I have tried to explore and even to celebrate in Metamagical Themas. (It is not for nothing that the word "magic" appears inside the title!) I hope that Metamagical Themas will help people to bring more clarity, precision, and elegance to their thinking about situations large and small. I also hope that it will inspire people to dedicate more of their energies to global problems in this lunatic but lovable world, because we live in a time of unprecedented urgency. If we do not care enough now, future generations may not exist to thank us for their existence and for our caring. 



Introduction                                                                      xxviii \sectiontitle {Section I:}

Snags and Snarls 



\sectiontitle {Section I:}

Snags and Snarls The title of this section conveys the image of problematical twistiness, The twists dealt with here are those whereby a system (sentence, picture, language, organism, society, government, mathematical structure, computer program, etc.) twists back on itself and closes a loop. A very general name for this is reflexivity. When realized in different ways, this abstraction becomes a concrete phenomenon. Examples are: self- reference, self-description, self-documentation, self-contradiction, self-questioning, self- response, self-justification, self-refutation, self-parody, self-doubt, self-definition, self- creation, self-replication, self-modification, self-amendment, self-limitation, self- extension, self-application, self-scheduling,-self-watching, and on and on. In the following four chapters, -these strange phenomena are illustrated in sentences and stories that talk about others, ideas that propagate themselves from mind to mind, machines that replicate themselves, and games that modify their own rules. The variety of these loopy tangles is quite remarkable, and the subject is far, far from being exhausted. Furthermore, although their connection with paradox may make reflexive systems seem no more than fin al playthings, study of them is of great importance in understanding many mathematical and scientific developments of this century, and is becoming ever more central to theories of intelligence and consciousness, whether natural or artificial Reflexivity will therefore make many return appearances in this book. 



1. On Self-Referential Sentences January, 1981 
I never expected to be writing a column for Scientific American. I remember once, years ago, wishing I were in Martin Gardner's shoes. It seemed exciting to be able to plunge into almost any topic one liked and to say amusing and instructive things about it to a large, well-educated, and receptive audience. The notion of doing such a thing seemed ideal, even dreamlike. Over the next several years, by a series of total coincidences (which turned out to be not so total), I met one after another of Martin's friends. First it was Ray Hyman, a psychologist who studies deception. He introduced me to the magician Jerry Andrus. Then I met the statistician and magician Persi Diaconis and the computer wizard Bill Gosper. Then came Scott Kim, and soon afterward, the mathematician Benoit Mandelbrot. All of a sudden, the world seemed to be orbiting Martin Gardner. He was at the hub of a magic circle, people with exciting, novel, often offbeat ideas, people with many-dimensional imaginations. Sometimes I felt overawed by the whole remarkable bunch. One day, five or so years ago, I had the pleasure of spending several hours with Martin in his house, discussing many topics, mathematical and otherwise. It was an enlightening experience for me, and it gave me a new view into the mind of someone who had contributed so much to my own mathematical education. Perhaps the most striking thing about Martin to me was his natural simplicity. I had been told that he is an adroit magician. This I found hard to believe, because one does not usually imagine someone so straightforward pulling the wool over anyone's eyes. However, I did not see him do any magic tricks. I simply saw his vast knowledge and love of ideas spread out before me, without the slightest trace of pride or pretense. The Gardners-Martin and his wife Charlotte-entertained me for the day. We ate lunch in the kitchen of their cozy three- story house. It pleased me somehow to see that there was practically no trace of mathematics or games or tricks in their simple but charming living room. After lunch-sandwiches that Martin and I made while standing by the kitchen sink-we climbed the two flights of stairs to Martin's hideaway. With his old typewriter and all kinds of curious jottings in an ancient filing cabinet 



and his legendary library of three-by-five cards, he reminded me of an old-time journalist, not of the center of a constellation of mathematical eccentrics and game addicts, to say nothing of magicians, anti-occultists, and of course the thousands of readers of his column. Occasionally we were interrupted by the tinkling of a bell attached to a string that led down the stairs to the kitchen, where Charlotte could pull it to get his attention. A couple of phone calls came, one from the logician and magician Raymond Smullyan, someone whose name I had known for a long time, but who I had no idea belonged to this charmed circle. Smullyan was calling to chat about a book he was writing on Taoism, of all things! For a logician to be writing about what seemed to me to be the most anti- logical of human activities sounded wonderfully paradoxical. (In fact, his book The Tao Is Silent is delightful and remarkable.) All in all, it was a most enjoyable day. Martin's act will be a hard one to follow. But I will not be trying to be another Martin Gardner. I have my own interests, and they are different from Martin's, although we have much in common. To express my debt to Martin and to symbolize the heritage of his column, I have kept his title "Mathematical Games" in the form of an anagram: "Metamagical Themas". What does "metamagical" mean? To me, it means "going one level beyond magic". There is an ambiguity here: on the one hand, the word might mean "ultramagical"-magic of a higher order-yet on the other hand, the magical thing about magic is that what lies behind it is always nonmagical. That's metamagic for you! It reflects the familiar but powerful adage "Truth is stranger than fiction." So my "Metamagical Themas" will, in Gardnerian fashion, attempt to show that magic often lurks where few suspect it, and, by the opposite token, that magic seldom lurks where many suspect it. 
\separator
In his July, 1979 column, Martin wrote a very warm review of my book Gödel, Escher, Bach: an Eternal Golden Braid. He began the review with a short quotation from my book. If I had been asked to guess what single sentence he would quote, I would never have been able to predict his choice. He chose the sentence "This sentences no verb." It is a catchy sentence, I admit, but something about seeing it again bothered me. I remembered how I had written it one day a few years earlier, attempting to come up with a new variation on an old theme, but even at the time it had not seemed as striking as I had hoped it would. After seeing it chosen as the symbol of my book, I felt challenged. I said to myself that surely there must be much cleverer types of self-referential sentence. And so one day I wrote down quite a pile of self-referential sentences and showed them to friends, which began a mild craze among a small group of us. In this column, I will present a selection of what I consider to be the cream of that crop. 



Before going further, I should explain the term "self-reference". Self-reference is ubiquitous. It happens every time anyone says "I" or., me" or "word" or "speak" or "mouth". It happens every time a newspaper prints a story about reporters, every time someone writes a book about writing, designs a book about book design, makes a movie about movies, or writes an article about self-reference. Many systems have the capability to represent or refer to themselves somehow, to designate themselves (or elements of themselves) within the system of their own symbolism. Whenever this happens, it is an instance of self-reference. Self-reference is often erroneously taken to be synonymous with paradox. This notion probably stems from the most famous example of a self-referential sentence, the Epimenides paradox. Epimenides the Cretan said, "All Cretans are liars." I suppose no one today knows whether he said it in ignorance of its self-undermining quality or for that very reason. In any case, two of its relatives, the sentences "I am lying" and "This sentence is false", have come to be known as the Epimenides paradox or the liar paradox. Both sentences are absolutely self-destructive little gems and have given self-reference a bad name down through the centuries. When people speak of the evils of self-reference, they are certainly overlooking the fact that not every use of the pronoun "I" leads to paradox. 
\separator
Let us use the Epimenides paradox as our jumping-off point into this fascinating land. There are many variations on the theme of a sentence that somehow undermines itself. Consider these two: 
This sentence claims to be an Epimenides paradox, but it is lying. This sentence contradicts itself-or rather-well, no, actually it doesn't! 
What should you do when told, "Disobey this command"? In the following sentence, the Epimenides quality jumps out only after a moment of thought: "This sentence contains exactly threee erors." There is a delightful backlash effect here. Kurt Gödel’s famous Incompleteness Theorem in metamathematics can be thought of as arising from his attempt to replicate as closely as possible the liar paradox in purely mathematical terms. With marvelous ingenuity, he was able to show that in any mathematically powerful axiomatic system S it is possible to express a close cousin to the liar paradox, namely, "This formula is unprovable within axiomatic system S. " In actuality, the Gödel construction yields a mathematical formula, not an English sentence; I have translated the formula back into English to show what he concocted. However, astute readers may have noticed that, strictly speaking, the phrase "this formula" has no referent, since when a formula 



is translated into an English sentence, that sentence is no longer a formula! If one pursues this idea, one finds that it leads into a vast space. Hence the following brief digression on the preservation of self-reference across language boundaries. How should one translate the French sentence Cette phrase en francais est difficile a traduire en anglais ? Even if you do not know French, you will see the problem by reading a literal translation: "This sentence in French is difficult to translate into English." The problem is: To what does the subject ("This sentence in French") refer? If it refers to the sentence it is part of (which is not in French), then the subject is self- contradictory, making the sentence false (whereas the French original was true and harmless); but if it refers to the French sentence, then the meaning of "this" is strained. Either way, something disquieting has happened, and I should point out that it would be just as disquieting, although in a different way, to translate it as: "This sentence in English is difficult to translate into French." Surely you have seen Hollywood movies set in France, in which all the dialogue, except for an occasional Bonjour or similar phase, is in English. What happens when Cardinal Richelieu wants to congratulate the German baron for his excellent command of French? I suppose the most elegant solution is -for him to say, "You have an excellent command of our language, mon cher baron ", and leave it at that. But let us undigress and return to the Gödelian formula and focus on its meaning. Notice that the concept of falsity (in the liar paradox) has been replaced by the more rigorously understood concept of provability. The logician Alfred Tarski pointed out that it is in principle impossible to translate the liar paradox exactly into any rigorous mathematical language, because if it were possible, mathematics would contain a genuine paradox -a statement both true and false-and would come tumbling down. Gödel’s statement, on the other hand, is not paradoxical, though it constitutes a hair-raisingly close approach to paradox. It turns out to be true, and for this reason, it is unprovable in the given axiomatic system. The revelation of Gödel’s work is that in any mathematically powerful and consistent axiomatic system, an endless series of true but unprovable formulas can be constructed by the technique of self-reference, revealing that somehow the full power of human mathematical reasoning eludes capture in the cage of rigor. In a discussion of Gödel’s proof, the philosopher Willard Van Orman Quine invented the following way of explaining how self-reference could be achieved in the rather sparse formal language Gödel was employing. Quine's construction yields a new way of expressing the liar paradox. It is this: 



"yields falsehood, when appended to its quotation." yields falsehood, when appended to its quotation. 
This sentence describes a way of constructing a certain typographical entity -namely, a phrase appended to a copy of itself in quotes. When you carry out the construction, however, you see that the end product is the sentence itself-or a perfect copy of it. (There is a resemblance here to the way self-replication is carried out in the living cell.) The sentence asserts the falsity of the constructed typographical entity, namely itself (or an indistinguishable copy of itself). Thus we have a less compact but more explicit version of the Epimenides paradox. It seems that all paradoxes involve, in one way or another, self-reference, whether it is achieved directly or indirectly. And since the credit for the discovery-or creation-of self-reference goes to Epimenides the Cretan, we might say: "Behind every successful paradox there lies a Cretan." On the basis of Quine's clever construction we can create a self-referential question: 
What is it like to be asked, "What is it like to be asked, self-embedded in quotes after its comma?" self-embedded in quotes after its comma? 
Here again, you are invited to construct a typographical entity that turns out, when the appropriate operations have been performed, to be identical with the set of instructions. This self-referential question suggests the following puzzle: What is a question that can serve as its own answer? Readers might enjoy looking for various solutions to it. 
\separator
When a word is used to refer to something, it is said to be being used. When a word is quoted, though, so that one is examining it for its surface aspects (typographical, phonetic, etc.), it is said to be being mentioned. The following sentences are based on this famous use-mention distinction: 
You can't have your use and mention it too. 
You can't have "your cake" and spell it "too". 
"Playing with the use-mention distinction" isn't "everything in life, you know". 
In order to make sense of "this sentence", you will have to ignore the quotes in "it". 



T'his is a sentence with "onions", "lettuce", "tomato", and "a side of fries to go". 
This is a hamburger with vowels, consonants, commas, and a period at the end. 
The last two are humorous flip sides of the same idea. Here are two rather extreme examples of self-referential use-mention play: 
Let us make a new convention: that anything enclosed in triple quotes-for example, "`No, I have decided to change my mind; when the triple quotes close, just skip directly to the period and ignore everything up to it"'-is not even to be read (much less paid attention to or obeyed). 

A ceux qui ne comprennent pas l'anglais, la phrase citee ci-dessous ne dit rien: "For those who know no French, the French sentence that introduced this quoted sentence has no meaning." 
The bilingual example may be more effective if you know only one of the_ two languages involved. Finally, consider this use-mention anomaly: "i should begin with a capital letter." This is a sentence referring to itself by the pronoun "I", a bit mauled, instead of through a pointing-phrase such as "this sentence"; such a sentence would seem to be arrogantly proclaiming itself to be an animate agent. Another example would be "I am not the person who wrote me." Notice how easily we understand this curious nonstandard use of "I". It seems quite natural to read the sentence this way, even though in nearly all situations we have learned to unconsciously create a mental model of some person-the sentence's speaker or writer-to whom we attribute a desire to communicate some idea. Here we take the "I" in a new way. How come? What kinds of cues in a sentence make us recognize that when the word "I" appears, we are supposed to think not about the author of the sentence but about the sentence itself? 
\separator
Many simplified treatments of Gödel’s work give as the English translation of his famous formula the following: "I am not provable in axiomatic system S. " The self-reference that is accomplished with such sly trickery in the formal system is finessed into the deceptively simple English word "I", and we can-in fact, we automatically do-take the sentence to be talking about itself. Yet it is hard for us to hear the following sentence as talking about itself: "I already took the garbage out, honey." The ambiguous referring possibilities of the first-person pronoun are a source of many interesting self-referential sentences. Consider these: . 



I am not the subject of this sentence. 
I am jealous of the first word in this sentence. 
Well, how about that-this sentence is about me! 
I am simultaneously writing and being written. 
This raises a whole new set of possibilities. Couldn't "I" stand for the writing instrument ("I am not a pen"), the language ("I come from Indo-European roots"), the paper ("Cut me out, twist me, and glue me to form a Mobius strip, please")? One of the most involved possibilities is that "I" stands not for the physical tokens we perceive before us but for some more ethereal and intangible essence, perhaps the meaning of the sentence. But then, what is meaning? The next examples explore that idea: 
I am the meaning of this sentence. 
I am the thought you are now thinking. 
I am thinking about myself right now. 
I am the set of neural firings taking place in your brain as you read the set of letters in this sentence and think about me. 
This inert sentence is my body, but my soul is alive, dancing in the sparks of your brain. 
The philosophical problem of the connections among Platonic ideas, mental activity, physiological brain activity, and the external symbols that trigger them is vividly raised by these disturbing sentences. This issue is highlighted in the self-referential question, "Do you think anybody has ever had precisely this thought before?" To answer the question, one would have to know whether or not two different brains can ever have precisely the same thought (as two different computers can run precisely the same program). An illustration of this possibility may be found in Figure 24-2. 1 have often wondered: Can one brain have the same thought more than once? Is a thought something Platonic, something whose essence exists independently of the brain it is occurring in? If the answer is "Yes, thoughts are brain-independent", then the answer to the self-referential question would also be yes. If it is not, then no one could ever have had the same thought before-not even the person thinking it! Certain self-referential sentences involve a curious kind of communication between the sentence and its human friends: 



You are under my control because I am choosing exactly what words you are made out of, and in what order. 
No, you are under my control because you will read until you have reached the end of me. 
Hey, down there-are you the sentence I am writing, or the sentence I am reading? 
And you up there-are you the person writing me, or the person reading me? 
You and I, alas, can have only one-way communication, for you are a person and , a mere sentence. 
As long as you are not reading me, the fourth word of this sentence has no referent. 
The reader of this sentence exists only while reading me. 
Now that is a rather frightening thought! And yet, by its own peculiar logic, it is certainly true. 
Hey, out there-is that you reading me, or is it someone else? 
Say, haven't you written me somewhere else before? 
Say, haven't I written you somewhere else before? 
The first of the three sentences above addresses its reader; the second addresses its author. In the last one, an author addresses a sentence. Many sentences include words whose referents are hard to figure out because of their ambiguity-possibly accidental, possibly deliberate: 
Thit sentence is not self-referential because "thit" is not a word. 
No language can express every thought unambiguously, least of all this one. 
In the Escher-inspired Figure 1-1, visual and verbal ambiguity are simultaneously exploited. 
\separator




FIGURE 1-1. Ambiguity: What is being described-the hand, or the writing? [Drawing by David Moser, after AL C. Escher. ] 
Let us turn to a most interesting category, namely sentences that deal with the languages they are in, once were in, or might have been in: 
When you are not looking at it, this sentence is in Spanish. 
I had to translate this sentence into English because I could not read the original Sanskrit. 
The sentence now before your eyes spent a month in Hungarian last year and was only recently translated back into English. 
If this sentence were in Chinese, it would say something else. 
.siht ekil ti gnidaer eb d'uoy werbeH ni erew ecnetnes siht fl 
The last two sentences are examples of counterfactual conditionals. Such a sentence postulates in its first clause (the antecedent) some contrary-to-fact situation (sometimes called a "possible world") and extrapolates in its second clause (the consequent) some consequence of it. This type of sentence opens up a rich domain for self-reference. Some of the more intriguing self-referential counterfactual conditionals I have seen are the following: 



If this sentence didn't exist, somebody would have invented it. If I had finished this sentence, 
If there were no counterfactuals, this sentence would not be paradoxical. 
If wishes were horses, the antecedent of this conditional would be true. 
If this sentence were false, beggars would ride. 
What would this sentence be like if it were not self-referential? 
What would this sentence be like if it were 3? 
Let us ponder the last of these (invented by Scott Kim) for a moment. In a world where π actually did have the value 3, you wouldn't ask about how things would be if it were 3. Instead, you might muse "if π were 2" or "if π weren't 3". So one's first answer to the question might be this: "What would this sentence be like if π weren't 3?". But there is a problem. The referent of "this sentence" has now changed identity. So is it fair to say that the second sentence is an answer to the first? It is a little like a woman who muses, "What would I be doing now if I had had different genes?" The problem is that she would not be herself; she would be someone else, perhaps the little boy across the street, playing in his sandbox. Personal pronouns like "I" cannot quite keep up with such strange hypothetical world-shifts. But getting back to Scott Kim's counterfactual, I should point out that there is an even more serious problem with it than so far mentioned. Changing the value of π is, to put it mildly, a radical change in mathematics, and presumably you cannot change mathematics radically without having radically changed the fabric of the universe within which we live. So it is quite doubtful that any of the concepts in the sentence would make any sense if π were 3 (including the concepts of "π", "3", and so on). 
Here are two more counterfactual conditionals to put in your pipe and smoke: 
If the subjunctive was no longer used in English, this sentence would be grammatical. 
This sentence would be seven words long if it were six words shorter. 
These two lovely examples, invented by Ann Trail (who is also responsible for quite a few others in this column), bring us around to sentences that comment on their own form. Such sentences are quite distinct from ones 



that comment on their own content (such as the liar paradox, or the sentence that says "This sentence is not about itself, but about whether it is about itself."). It is easy to make up a sentence that refers to its own form, but it is hard to make up an interesting one. Here are a few more quite good ones: 
because I didn't think of a good beginning for it. 
This sentence was in the past tense. 
This sentence has contains two verbs. 
This sentence contains one numeral 2 many. 
a preposition. This sentence ends in 
In the time it takes you to read this sentence, eighty-six letters could have been processed by your brain. 
\separator
David Moser, a composer and writer, is a detector and creator of self-reference and frame-breaking of all kinds. He has even written a story in which every sentence is self- referential (it is included in Chapter 2). It might seem unlikely that in such a limited domain, individual styles could arise and flourish, but David has developed a self- referential style quite his own. As a mutual friend (or was it David himself?) wittily observed, "If David Moser had thought up this sentence, it would have been funnier." Many Moser creations have been used above. Some further Moserian delights are these: 
This is not a complete. Sentence. This either. 
This sentence contains only one nonstandard English flutzpah. 
This gubblick contains many nonsklarkish English flutzpahs, but the overall pluggandisp can be glorked from context. 
This sentence has cabbage six words. 
In my opinion, it took quite a bit of flutzpah to just throw in a random word so that there would be cabbage six words in the sentence. That idea inspired the following: "This sentence has five (5) words." A few more miscellaneous Moserian gems follow: 



This is to be or actually not two sentences to be, that is the question, combined 
It feels sooo good to have your eyes run over my curves and serifs. 
This sentence is a !!!! premature punctuator 
Sentences that talk about their own punctuation, as the preceding one does, can be quite amusing. Here are two more: 
This sentence, though not interrogative, nevertheless ends in a question mark? 
This sentence has no punctuation semicolon the others do period 
Another ingenious inventor of self-referential sentences is Donald Byrd, several of whose sentences have already been used above. Don too has his own very characteristic way of playing with self-reference. Two of his sentences follow: 
This hear sentence do'nt know Inglish purty good. 
If you meet this sentence on the board, erase it. 
The latter, via its form, alludes to the Buddhist saying "If you meet the Buddha on the road, kill him." Allusion through similarity of form is, I have discovered, a marvelously rich vein of self-reference, but unfortunately this article is too short to contain a full proof of that discovery. I shall explicitly discuss only two examples. The first is "This sentence verbs good, like a sentence should." Its primary allusion is to the famous slogan "Winston tastes good, like a cigarette should", and its secondary allusion is to, "This sentence no verb." The other example involves the following lovely self-referential remark, once made by the composer John Cage: "I have nothing to say, and I am saying it." This allows the following rather subtle twist to be made: "I have nothing to allude to, and I am alluding to it." 
\separator
Some of the best self-referential sentences are short but sweet, relying for their effect on secondary interpretations of idiomatic expressions or well-known catch phrases. Here are five of my favorites, which seem to defy other types of categorization: 
Do you read me? 



This point is well taken. 
You may quote me. 
I am going two-level with you. 
I have been sentenced to death. 
In some of these, even sophisticated non-native speakers would very likely miss what's going on. Surely no article on self-reference would be complete without including a few good examples of self-fulfilling prophecy. Here are a few: 
This prophecy will come true. 
This sentence will end before you can say `Jack Rob 
Surely no article on self-reference would be complete without including a few good examples of self-fulfilling prophecy. 
Does this sentence remind you of Agatha Christie? 
That last sentence-one of Ann Trail's-is intriguing. Clearly it has nothing to do with Agatha Christie, nor is it in her style, and so the answer ought to be no. Yet I'll be darned if I can read it without being reminded of Agatha Christie! (And what is even stranger is that I don't know the first thing about Agatha Christie!) In closing, I cannot resist the touching plea of the following Byrdian sentence: 
Please, oh please, publish me in your collection of self-referential sentences! 

Post Scriptum. This first column of mine triggered a big wave of correspondence, some of which is presented in the next chapter. Most of the correspondence was light-hearted, but there were a number of serious letters that intrigued me. Here is a repartee that appeared in the pages of Scientific American a few months later. 
The kind of structural analysis engaged in, and the resulting questions raised by, Douglas Hofstadter in his amusing and intriguing article concerning self- referential sentences need not lead inevitably to bafflement of the reader. 



Help is at hand from the "laggard science" psychology, but only from that carefully defined quarter of psychology known as behavior analysis, which was pt ogenerated by the famous Harvard psychologist B. F. Skinner almost 50 years ago. In examining the implications of linguistic analyses such as Hofstadter's for the serious student of verbal behavior, Skinner comments in his book About Behaviorism (pages 98-99) as follows: 
Perhaps there is no harm in playing with sentences in this way or in analyzing the kinds of transformations which do or do not make sentences acceptable to the ordinary reader, but it is still a waste of time, particularly when the sentences thus generated could not have been emitted as verbal behavior. A classical example is a paradox, such as 'This sentence is false', which appears to be true if false and false if true. The important thing to consider is that no one could ever have emitted the sentence as verbal behavior. A sentence must be in existence before a speaker can say, 'This sentence is false', and the response itself will not serve, since it did not exist until it was emitted. What the logician or linguist calls a sentence is not necessarily verbal behavior in any sense which calls for a behavioral analysis. 
As Skinner pointed out long ago, verbal behavior results from contingencies of reinforcement arranged by verbal communities, and it is these contingencies that must be analyzed if we are to identify the variables that control verbal behavior. Until we grasp the full import of Skinner's position, which goes beyond structure to answer why we behave as we do verbally or nonverbally, we shall continue to fall back on prescientific formulations that are about as useful in understanding these phenomena as Hofstadter's quaint metaphorical speculation: "Such a sentence would seem to be arrogantly proclaiming itself to be an animate agent." George Brabner College of Education University of Delaware 
I felt compelled to reply to Professor Brabner's interesting views about these matters, and so here is what I wrote: 
I assume that the quote from B. F. Skinner reflects Professor Brabner's own sentiments about the likelihood of self-referential utterances. I am always baffled by people who doubt the likelihood of self-reference and paradox. Verbal behavior comes in many flavors. Humor, particularly self-referential humor, is one of the most pervasive flavors of verbal behavior in this century. One has only to watch the Muppets or Monty Python on television to see dense and intricate webs of self- reference. Even advertisements excel in self-reference. In art, Rene Magritte, Pablo Picasso, M. C. Escher, John Cage, and dozens of others have played with the level-distinction between that which represents and that which is represented. The "artistic behavior" that results includes much self-reference and many confusing and sometimes exhilaratingly paradoxical 


tangles. Would Professor Brabner say that no one could ever have "emitted" such works as "artistic behavior"? Where is the borderline? Ordinary language, as I pointed out in my column, is filled with self-reference, usually a little milder-seeming than the very sharply pointed paradoxes that Professor Brabner objects to. "Mouth", "word", and so on are all self-referential. Language is inherently filled with the potential of sharp turns on which it may snag itself. Many scholarly papers begin with a sentence about "the purpose of this paper". Newspapers report on their own activities, conceivably on their own inaccuracies. People say, "I'm tired of this conversation." Arguments evolve about arguments, and can get confusingly and painfully self-involved. Has Professor Brabner never thought of "verbal behavior" in this light? It is likely that in hunting woolly mammoths, no one found it extraordinarily likely to shout, "This sentence is false!" However, civilization has come a long way since those days, and the primitive purposes of language have by now been almost buried under an avalanche of more complex purposes. Part of human nature is to be introspective, to probe. Part of our "verbal behavior" deliberately, often playfully, explores the boundaries between conceptual levels of systems. All of this has its root in the struggle to survive, in the fact that our brains have become so flexible that much of their time is spent in dealing with their own activities, consciously or unconsciously. It is simply a consequence of representational power-as Kurt Godel showed-that systems of increasing complexity become increasingly self-referential. It is quite possible for people filled with self-doubt to recognize this trait in themselves, and to begin to doubt their self-doubt itself. Such psychological dilemmas are at the heart of some current theories of therapy. Gregory Bateson's "double bind", Victor Frankl's "logotherapy", and Paul Watzlawick's therapeutic ideas are all based on level-crossing paradoxes that crop up in real life. Indeed, psychotherapy is itself based completely on the idea of a "twisted system of self--a self that wants to reach inward and change some presumably wrong part of itself. We human beings are the only species to have evolved humor, art, language, tangled psychological problems, even an awareness of our own mortality. Self- reference-even of the sharp Epimenides type-is connected to profound aspects of life. Would Professor Brabner argue that suicide is not conceivable human behavior? Finally, just suppose Professors Skinner and Brabner are right, and no one ever says exactly "This sentence is false." Would this mean that study of such sentences is a waste of time? Still not. Physicists study ideal gases because they represent a distillation of the most significant principles of the behavior of real gases. Similarly, the Epimenides paradox is an "ideal paradox"-one that cuts crisply to the heart of the matter. It has opened up vast domains in logic, pure science, philosophy, and other disciplines, and will continue to do so despite the skepticism of behaviorists. 
It is a curious coincidence that the only other reply to my article that was printed in the "Letters" column of Scientific American also came from the University of Delaware. Here it is: 



I hope that you do not receive any correspondence concerning Douglas R. Hofstadter's article on self-reference. I should like to inform your readers that many years of study on this problem have convinced me no conclusion whatsoever can be drawn from it that would stand up to a moment's scrutiny. There is no excuse for Scientific American to publish letters from those cranks who consider such matters to be worthy of even the slightest notice. 
A. J. Dale Department of Philosophy University of Delaware 
I replied as follows: 
Many years of reading such letters have convinced me that no reply whatsoever can be given to them that would stand up to a moment's scrutiny. There is no excuse for publishing responses to those cranks who send them. 
After these two exchanges had appeared in print, a number of people remarked to me that they'd read the two letters from Delaware that had attacked me, and had enjoyed my responses. Two; I guess it wasn't so obvious that Dale's letter was completely tongue-in- cheek. In fact, that was its point. 
\separator
Two other letters stand out sharply in my memory. One was from an individual who signed himself (I presume it is a male) as "Mr Flash gFiasco". Mr Flash insisted that a sentence cannot say what it shows. The former concerns only its content, which is supposedly independent of how it manifests itself in print, while the latter is a property exclusively of its form, that is, of the physical sentence only when it is in print. This distinction sounds crystal-clear at first, but in reality it is mud-blurry. Here is some of what Flash wrote me: 
For a sentence to attempt to say what it shows is to commit an error of logical types. It seems to be putting a round peg into a square hole, whereas it is instead putting a round peg into something which is not a hole at all, square or otherwise. This is a category mismatch, not a paradox. It is like throwing the recipe in with the flour and butter and eggs. The source of the equivocation is an illegitimate use of the term 'this'. 'This' can point to virtually anything, but 'this' cannot point to itself. If you stick out your index finger, you can point to virtually anything; and by curling it you can even point to the pointing finger; but you cannot point to pointing. Pointing is of a higher logical type than the thing which is doing the pointing. Similarly, the referent of 'this sentence' can be virtually anything but that sentence. Sentences of the form exemplified by 'This sentence no verb.' and 'This sentence has a verb.' are not well- formed: they commit fallacies of logical type equivocation. Thus their self-referential character is not genuine and they present no problem as paradoxes. 

There will always be people around who will object in this manner, and in the Brabnerian manner. Such people think it is possible to draw a sharp line between attributes of a printed sentence that can be considered part of its form (e.g., the typeface it is printed in, the number of words it contains, and so on), and attributes that can be considered part of its content (i.e., the things and events and relationships that it refers to). Now, I am used to thinking about language in terms of how to get a machine to deal with it, since I look at the human brain as a very complex machine that can handle language (and many other things as well). Machines, in trying to make sense of sentences, have access to nothing more than the form of such sentences. The content, if it is to be accessible to a machine, has to be derived, extracted, constructed, or created somehow from the sentence's physical structure, together with other knowledge and programs already available to the machine. When very simple processing is used to operate on a sentence, it is convenient to label the information thus obtained "syntactic". For instance, it is clearly a syntactic fact about "This sentence no verb." that it contains six vowels. The vowel-consonant distinction is obviously a typographical one, and typographical facts are considered superficial and syntactic. But there is a problem here. With different depths of processing, aspects of different degrees of "semanticity" may be detected. Consider, for example, the sentence "Mary was sick yesterday." Let's call it Sentence M. Listed below are the results of seven different degrees of processing of Sentence M by a hypothetical machine, using increasingly sophisticated programs and increasingly large knowledge bases. You should think of them as being English translations, for your convenience, of computational structures inside the machine that it can act on and use fluently. 
1. Sentence M contains twenty characters. 2. Sentence M contains four English words. 3. Sentence M contains one proper noun, one one adverb, in that order. . 4. Sentence,M contains one human's name, one linking verb, one adjective describing a potential health state of a living being, and one temporal adverb, in that order. 5. The subject of Sentence M is a pointer to an individual named `Mary', the predicate is an ascription of ill health to the individual so indicated, on the day preceding the statement's utterance. 6. Sentence M asserts that the health of an individual named 'Mary' was not good the day before today. 7. Sentence M says that Mary was sick yesterday. verb, one adjective, and 
Just where is the boundary line that says, "You can't do that much processing!"? A machine that could go as far as version 7 would have 



actually understood-at least in some rudimentary sense-the content of Sentence M. Work by artificial-intelligence researchers in the field of natural language understanding has produced some very impressive results along these lines, considerably more sophisticated than what is shown here. Stories can be "read" and "understood", at least to the extent that certain kinds of questions can be answered by the machine when it is probed for its understanding. Such questions can involve information not explicitly in the story itself, and yet the machine can fill in the missing information and answer the question. I am making this seeming digression on the processing of language by computers because intelligent people like Mr Flash qFiasco seem to have failed to recognize that the boundary line between form and content is as blurry as that between blue and green, or between human and ape. This comparison is not made lightly. Humans are supposedly able to get at the "content" of utterances, being genuine language-users, while apes are not. But ape-language research clearly shows that there is some kind of in-between world, where a certain degree of content can be retrieved by a being with reduced mental capacity. If mental capacity is equated with potential processing depth, then it is obvious why it makes no sense to draw an arbitrary boundary line between the form and the content of a sentence. Form blurs into content as processing depth increases. Or, as I have always liked to say, "Content is just fancy form." By this I mean, of course, that "content" is just a shorthand way of saying "form as perceived by a very fancy apparatus capable of making complex and subtle distinctions and abstractions and connections to prior concepts". Flash qFiasco's down-home, commonsense distinction between form and content breaks down swiftly, when analyzed. His charming image of someone making a "category error" by throwing a recipe in with the flour and butter and eggs reveals that he has never had Recipe Cake. This is a delicious cake whose batter is made out of cake recipes (if you use pie recipes, it won't taste nearly as good). The best results are had if the recipes are printed in French, in Baskerville Roman. A preponderance of accents aigus lends a deliciously piquant aroma to the cake. My recommendation to Brabner and qFiasco is: "Let them eat recipes." 
\separator
Finally, I come to John Case, a computer scientist who wrote from Yale, insisting that there is no conceptual problem whatsoever in translating the French sentence "Cette phrase en franfais est dicile a traduire en anglais " into English. Case's translation was the following English sentence: 
The French sentence "Cette phrase en franfais est difficile a traduire en anglais" is difficult to translate into English. 



In other words, Case translates a self-referential French sentence into an other- referential English sentence. The English sentence talks about the French sentence-in fact it quotes it completely! Something radical is missing here. At one level, of course, Case is right: now the two sentences, one French and one English, both are talking about (or pointing to) the same thing (the French sentence). But the absolute crux of the French one is its tangledness; the English one completely lacks that quality. Clearly Case has had to make a sacrifice, a compromise. The alternative, which I prefer, is to construct in English an analogue to the French sentence: a self-referential English sentence, one that has a tangledness isomorphic to that of the French sentence. That's where the essence of the sentence lies, after all! "But is that its translation ?" you might ask. A good question. lonesco once remarked, "The French for London is Paris." (Use-mention fanatic that I am, I assume that he meant "The French for `London' is `Paris' ", although it is pungent either way.) What he meant was that in understanding situations, French people tend to translate them into their own frame of reference. This is of course true for all of us. If Mary tells Ann, "My brother died", and if Ann does not know Mary's brother, then how can she understand this statement? Surely projection is of the essence: Ann will imagine her own brother dying (if she has one-and if not, then her sister, a good friend, possibly even a pet!). This alternate frame of reference allows Ann to empathize with Mary. Now if Ann did know Mary's brother somewhat, then she might flicker between thinking of him as the person she vaguely remembers and thinking of her own brother (friend, pet, or whatever) dying. This dilemma (discussed further in the postscript to \chapter {24}
into what they would be for me, or do I stand apart and survey them completely objectively and impassively? Case is advocating the latter, which is all very well as an intellectual stance to adopt, but when it comes to real life, it just won't cut the mustard. To be concrete, one might ask: What was the actual solution used in the French edition of Scientific American ? The answer, surprising no one, I hope, was this: "This English sentence is difficult to translate into French." I rest my case. 
\separator
I wonder what literalists like John Case would suggest as the proper translation of the title of the book All the President's Men (a book about the downfall of President Nixon, a downfall that none of the people around him could prevent). Would they say that Tous les hommes du President fills the bill admirably? Back-translated rather literally, it means "All the men of the President". It completely lacks the allusion -the reference by similarity of form-to the nursery rhyme "Humpty Dumpty". Is that dispensable? In my 



opinion, hardly. To me, the essence of the title resides in that allusion. To lose that allusion is to deflate the title totally. Of course, what do I mean by "that allusion"? Do I wish the French title to contain, somehow, an allusion to an English nursery rhyme? That would be rather pointless. Well, then, do I want the French title to allude to the French version of "Humpty Dumpty"? It all depends how well known that is. But given that Humpty Dumpty is practically an unknown figure to French-speaking people, it seems that something else is wanted. Any old French nursery rhyme? Obviously not. The critical allusion is to the lines "All the King's horses/ And all the King's men/ Couldn't put Humpty together again." Are there-anywhere in French literature-lines with a similar import? If not, how about in French popular songs? In French proverbs? Fairy tales? One might well ask why French-speaking people would ever care about reading a book about Watergate in the first place. And even if they did want to read it, shouldn't it be completely translated, so that it happens in a French-speaking city? Come to think of it, didn't loratno once remark that the French for Washington is Montreal? Clearly, this is carrying things to an extreme. There must be some middle ground of reasonableness. These are matters of subtle judgment, and they are where being human and flexible makes all the difference. Rigid rules about translation may lead you to a kind of mechanical consistency, but at the sacrifice of all depth and charm. The problem of self-referential sentences is just the tip of the iceberg, as far as translation is concerned. It is just that these issues show up very early when direct self-reference is concerned. When self-reference (or reference in general, for that matter) is indirect, mediated by form, then fluidity is required. The understanding of such sentences involves a mixture of deriving the content and yet retaining the form in mind, letting qualities of the form conjure up flavors and enhance the meaning with a halo of not-quite-conscious pseudo-meanings, connotations, flavors, that flicker in the mind, not quite in reach, not quite out of reach. Self-reference is a good starting point for investigation of this kind of issue, because it is so much on the surface there. You can't sweep the problems under the rug, even though some would like to do so. 
\separator
This first column, together with this postscript, provides a good introduction to the book as a whole, because many central issues are touched on: codes, translation, analogies, artificial intelligence, language and machines, mind and meanings, self and identity, form and content-all the issues I originally was motivated by when first writing that collection of teasing self-referential sentences. 



Self-Referential Sentences: A Follow-Up 
January, 1982 


AS January has rolled around again, I thought I'd give a follow-up to my column of a year ago on self-referential sentences, and that is what this column is; however, before we get any further, I would like to take advantage of this opening paragraph to warn those readers whose sensibilities are offended by explicit self-referential material that they probably will want to quit reading before they reach the end of this paragraph, or for that matter, this sentence-in fact, this clause-even this noun phrase-in short, this. Well, now that we've gotten that out of the way, I would like to say that, since last January, I have received piles upon piles of self-referential mail. Tony Durham astutely surmised: "What with the likely volume of replies, I should not think you are reading this in person." John C. Waugh's letter yelped: "Help, I'm buried under an avalanche of reader's responses!" At first, I thought Waugh himself was empathizing with my plight, putting words into my own mouth, but then I realized it was his letter calling for help. Fortunately, it was rescued, and now is comfortably nestled in a much reduced pile. Indeed, I have had to cull from that massive influx of hundreds of replies a very small number. Here I shall present some of my favorites. Before leaving the topic of mail, I would like to point out that the postmark on Ivan Vince's postcard from Britain cryptically remarked, "Be properly addressed." Was this an order issued by the post office to the postcard itself? If so, then British postcards must be far more intelligent than American ones; I have yet to meet a postcard that could read, let alone correct its own address. (One postcard that reached me was addressed to me in care of Omni magazine! And yet somehow it arrived.) I was flattered by a couple of self-undermining compliments. Richard Ruttan wrote, "I just can't tell you how much I enjoyed your first article.", and John Collins said, "This does not communicate my delight at January's column." I was also pleased to learn that my fame had spread as far as the men's room at the Tufts University Philosophy Department, where Dan 



Dennett discovered "This sentence is graffiti. -Douglas R. Hofstadter" penned on the wall. 
\separator

A popular pastime was the search for interesting self-answering questions. However, only a few succeeded in genuinely ' jootsing" (jumping out of the system), which, to me, means being truly novel. It seems that successes in this limited art form are not easy to come by. John Flagg cynically remarked (I paraphrase slightly): "Ask a self- answering question, and get a self-questioning answer." One of my favorites was given by Henry Taves: "I fondly remember a history exam I encountered in boarding school that contained the following: 'IV. Write a question suitable for a final exam in this course, and then answer it.' My response was simply to copy that sentence twice." I was delighted by this. Later, upon reflection, I began to suspect something was slightly wrong here. What do you think? Richard Showstack contributed two droll self-answering questions: "What question no verb?" and "What is a question that mentions the word `umbrella' for no apparent reason?" Jim Shiley sent in a clever entry that I modify slightly into "Is this a rhetorical question, or is this a rhetorical question?" He also contributed the following idea: 
Take a blank sheet of paper and on it write: 
How far across the page will this sentence run? 
Now if some polyglot friend of yours points out that the same string of phonemes in Ural-Altaic means '2.3 inches', send me a free subscription to Scientific American. Otherwise, if the inscription of a question counts both as the question and as unit of measure, I at least get a booby prize. But I think somehow I bent the rules. 
My own solutions to the problem of the self-answering question are actually not so much self-answering as self-provoking, as in the following example: "Why are you asking me that out of the blue?" It is obvious that when the question is asked out of the blue, it might well elicit an identical response, indicating the hearer's bewilderment. Philip Cohen relayed the following anecdote about a self-answering question, from Damon Knight: "Terry Carr, an old friend, sent us a riddle on a postcard, then the answer on another postcard. Then he sent us another riddle: `How do you keep a turkey in suspense?' and never sent the answer. After about two weeks, we realized that was the answer." 
\separator

Several of the real masterpieces sent in belong to what I call the self-documenting category, of which a simple example is Jonathan Post's "This 


sentence contains ten words, eighteen syllables and sixty-four letters." A neat twist is supplied by John Atkins in his sentence " 'Has eighteen letters' does." The self- documenting form can get much more convoluted and introspective. An example by the wordplay master Howard Bergerson was brought to my attention by Philip Cohen. It goes: 
In this sentence, the word and occurs twice, the word eight occurs twice, the word four occurs twice, the word fourteen occurs four times, the word in occurs twice, the word seven occurs twice, the word the occurs fourteen times, the word this occurs twice, the word times occurs seven times, the word twice occurs eight times and the word word occurs fourteen times. 
That is good, but the gold medal in the category is reserved for Lee Sallows, who submitted the following tour de force: 
Only the fool would take trouble to verify that his sentence was composed of ten a's, three b's, four c's, four d's, forty-six e's, sixteen f's, four g's, thirteen h's, fifteen i's, two k's, nine l's, four m's, twenty-five n's, twenty-four o's, five p's, sixteen r's, forty-one s's, thirty-seven t's, ten u's, eight v's, eight w's, four x's, eleven y's, twenty-seven commas, twenty-three apostrophes, seven hyphens, and, last but not least, a single ! 
I (perhaps the fool) did take trouble to verify the whole thing. First, though, I carried out some spot checks. And I must say that when the first random spot check worked (I think I checked the number of `g's), this had a strong psychological effect: all of a sudden, the credibility rating of the whole sentence shot way up for me. It strikes me as weird (and wonderful) how, in certain situations, the verification of a tiny percentage of a theory can serve to powerfully strengthen your belief in the full theory. And perhaps that's the whole point of the sentence! The noted logician Raphael Robinson submitted a playful puzzle in the self- documenting genre. Readers are asked to complete the following sentence: 
In this sentence, the number of occurrences of 0 is __, of 1 is __, of 2 is __, of 3 is __, of 4 is __, of 5 is _, of 6 is __, of 7 is __, of 8 is __, and of 9 is _. 
Each blank is to be filled with a numeral of one or more digits, written in decimal notation. Robinson states that there are exactly two solutions. Readers might also search for two sentences of this form that document each other, or even longer loops of that kind. Clearly the ultimate in self-documentation would be a sentence that does more than merely inventory its parts; it would be a sentence that includes a rule as well, telling all the King's men how to put those parts back together again to create a full sentence-in short, a self-reproducing sentence. Such 



a sentence is Willard Van Orman Quine's English rendition of Kurt Gödel’s classic metamathematical homage to Epimenides the Cretan: 
"yields falsehood when appended to its quotation." yields falsehood when appended to its quotation. 
Quine's sentence in effect tells the reader how to construct a replica of the sentence being read, and then (just for good measure) adds that the replica (not itself for heaven's sake!) asserts a falsity! It's a bit reminiscent of the famous remark made by Epilopsides the Concretan (second cousin of Epimenides) to Flora, a beautiful young woman whose ardent love he could not return (he was betrothed to her twin sister Fauna): "Take heart, my dear. I have a suggestion that may cheer you up. Just take one of these cells from my muscular biceps here, and clone it. You'll soon wind up with a dashing blade who looks and thinks just like me! But do watch out for him-he is given to telling beautiful women real whoppers!" 
\separator

In the early 1950's, John von Neumann worked hard trying to design a machine that could build a replica of itself out of raw materials. He came up with a theoretical design consisting of hundreds of thousands of parts. Seen in hindsight and with a considerable degree of abstraction, the idea behind von Neumann's self-reproducing machine turns out to be pretty similar to the means by which DNA replicates itself. And this in turn is close to Gödel’s method of constructing a self-referential sentence in a mathematical language in which at first there seems to be no way of referring to the language itself. The First Every-Other-Decade Von Neumann Challenge is thus hereby presented for ambitious readers: Create a comprehensible and not unreasonably long self- documenting sentence that not only lists its parts (at the word level or, better yet, the letter level) but also tells how to put them together so that the sentence reconstitutes itself. (Notice, by the way, the requirement is that the sentence be not unreasonably long, which is different -very different-from being reasonably long.) The parts list (or seed) should be an inventory of words or typographical symbols, more or less as in the sentences created by Howard Bergerson and Lee Sallows. The inventoried symbols should in some way be clearly distinguishable from the text that talks about them. For instance, they can be enclosed in quotation marks, printed in another typeface, or referred to by name. It is not so important what convention is adopted, so long as the distinction is sharp. The rest of the sentence (the building rule) should be printed normally, since it is to be regarded not as typographical raw material but as a set of instructions. This is the use-mention distinction I discussed in Chapter 1, and to disregard it 



is a serious conceptual weakness. (It is a flaw in Sallows' sentence that slightly tarnishes the gold on his medal.) The building rule may not talk about normally-printed material-only about parts of the inventory. Thus, it is not permitted for the building rule to refer to itself in any way! The building rule has to describe structure explicitly. Furthermore (and this is the subtlest and probably the most often overlooked aspect of self-reference), the building rule must specify which parts are to be printed normally and which parts in quotes (or however the raw materials are being indicated). In this respect, Bergerson's sentence fails. Although, to its credit, it sharply distinguishes between use and mention by relying on upper case for the names of inventory items and lower case for item counts and filler words, it does not have separate inventories for items in upper case and lower case. Instead it lumps the two together, blurring a vital distinction. In the Von Neumann Challenge, extra points will be awarded for solutions given in Basic English, or whose seed is entirely at the letter level (as in Sallows' sentence). The Quine sentence, although it clearly incorporates a seed (the seven-word phrase in quotation marks) and a building rule (that of appending something to its quotation), is not a legal entry because its seed is too far from being raw material. It is so structured that it is like a fetus more than it is like a zygote. 
\separator

There is a very good reason, by the way, that the Quine sentence's seed is so complicated-in fact, is identical with the building rule, except for the quotation marks. The reason is simple to state: You've got to build a copy of the building rule out of raw materials, and the more your building rule looks like your seed, the simpler it will be to build a copy of it from a copy of the seed. To make a full new sentence, all you need to do is make two copies of the seed, carry out whatever simple manipulations will convert one copy of the seed into the building rule, and then splice the other copy of the seed onto the newly minted building rule to make up a complete new sentence, fresh off the assembly line. To make this clearer, it is helpful to show a slight variation on Quine's sentence. Imagine that you could recognize only the lowercase roman letters, and that uppercase letters were alien to you. Then text printed in upper case would, for all practical purposes, be devoid of meaning or interest, whereas text in lower case would be full of meaning and interest, able to suggest ideas or actions in your mind. Now suppose someone gave you a conversion table that matched each uppercase letter with its lowercase counterpart, so that you could "decode" uppercase text. Then one day you came across this piece of "meaningless" uppercase text: 



YIELDS A FALSEHOOD WHEN USED AS THE SUBJECT OF ITS LOWERCASE VERSION 
On being decoded, it would yield a lowercase sentence, or rather, a lowercase sentence fragment-a predicate without a subject. Suggestive, eh? What might you try out, as a possible subject of that predicate? This notion of two parallel alphabets, one in which text is inert and meaningless and the other in which text is active and meaningful, may strike you as yielding no more than a minor variation on Quine's sentence, but in fact it is very similar to an exceedingly clever trick that nature discovered and has exploited in every cell of every living organism. Our seed-our genome-our DNA-is a huge long volume of inert text written in a chemical alphabet that has 64 "uppercase" letters (codons). Our building rules-our enzymes-are short, pithy slogans of active text written in a different chemical alphabet that has just twenty "lowercase" letters (amino acids). There is a map (the genetic code) that converts uppercase letters into lowercase ones. Obviously, some lowercase letters must correspond to more than one uppercase letter, but here that is a detail. It also turns out that three characters of the uppercase alphabet are not letters but punctuation marks telling where one pithy slogan ends and the next one begins-but again, these are details. \intro (See Chapter 27 for some of those details.) \endintro
Once you know this mapping, you often won't even remember to distinguish between the two chemical alphabets: the inert uppercase codon alphabet and the active lowercase amino acid alphabet. The main thing is that, armed with the genetic code, you can read the DNA book (seed) as if it were a sequence of enzyme slogans (building rules) telling how to write a new DNA book together with a new set of enzyme slogans! It is a perfect parallel to our variation on the Quine sentence, where inert, uppercase seed-text was, converted into active, lowercase rule-text that told how to make a copy of the full Quine sentence, given its seed. A cell's DNA and enzymes act like the seed and building rules of Quine's sentence, or the parts list and building rules of von Neumann's self-reproducing automaton-or then again, like the seed and building rules of computer programs that print themselves out. It is amazing how universal this mechanism of self-reference is, and for that reason I always find it quaint that people who rant and rave against the silliness of self-reference are themselves composed of trillions and trillions of tiny self-referential molecules. 
\separator

Scott Kim and I constructed an intriguing pair of sentences: 
The following sentence is totally identical with this one, except that the words `following' and `preceding' have been exchanged, as have the words 'except' and `in', and the phrases `identical with' and `different from'. 



The preceding sentence is totally different from this one, in that the words preceding' and `following' have been exchanged, as have the words `in' and except', and the phrases 'different from' and `identical with'. 
At first glance, these sentences are reminiscent of a two-step variant on the Epimenides paradox ("The following sentence is true."; "The preceding sentence is false."). On second glance, though, they are seen to say exactly the same thing. Curiously, my Australian colleague and sometime alter ego, Egbert B. Gebstadter, writing in his ever fascinating but often-furiating monthly row "Thetamagical Memas" (which appears in Literary Australian), disagrees with me; he maintains they say totally different things. (See figure 2-1.) Not surprisingly, several of the sentences submitted by readers had a paradoxical flavor. Some were variants on Bertrand Russell's paradox about the barber who shaves all those who do not shave themselves, or the set of all sets that do not include themselves as elements. For instance, Gerald Hull concocted this strange sentence: "This sentence refers to every sentence that does not refer to itself." Is Hull's concoction self-referential, or is it not? In a similar vein, Michael Gardner cited a Reed College senior thesis whose dedication ran: "This thesis is dedicated to all those who did not dedicate their theses to themselves." The book Model Theory, by C. C. Chang and H. J. Keisler, bears a similar dedication, as Charles Brenner pointed out to me. He also suggested another variant on Russell's paradox: Write a computer program that prints out a list of all programs that do not ever print themselves out. The question is, of course: Will this program ever print itself out? One of the most disorienting sentences came from Robert Boeninger: "This sentence does in fact not have the property it claims not to have." Got that? A serious problem seems to be to figure out just what property it is that the sentence claims it lacks. The Dutch mathematician Hans Freudenthal sent along a charming paradoxical anecdote based on self-reference: 
There is a story by the eighteenth-century German Christian Gellert called "Der Bauer and sein Sohn" ("The Peasant and His Son"). One day during a walk, when the son tells a big lie, his father direly warns him about the "Liars' Bridge", which they are approaching. This bridge always collapses when a liar walks across it. After hearing this frightening warning, the boy admits his lie and confesses the truth. When I [Freudenthal] told a ten-year-old boy this story, he asked me what happened when they eventually came to the bridge. I replied, "It collapsed under the father, who had lied, since in fact there is no Liars' Bridge." (Or did it?) 

C. W. Smith, writing from London, Ontario, described a situation reminiscent of the Epimenides paradox: 



Thetamagical Memas Seeking the Whence of Letter and Spirit EGBERT B. GEBSTADTER 



A Copious Concatenation of Artsy, Scientistic, and LiterL Mumbo-Jumbo FIGURE 2-1. The cover of Egbert B. Gebstadter's latest book, showing some of his " Whorly Art. " See the Bibliography for a short description of the book. Gebstadter, best known as the author of Copper, Silver, Gold: an Indestructible Metallic Alloy, also co-edited The Brain's U with Australian philosopher Denial E. Dunnitt, and for two and a ha f years wrote a monthly row ("Thetamagical Memas ")for Literary Australian. Having spent the last several years in the Psychology Department of Pakistania University in Willington, Pakistania, he has recently joined the faculty of the Computer Science Department of the University of Mishuggan in Tom Treeline, Mishuggan, where he occupies the Rexall Chair in the College ofArt, Sciences, and Letters. His current research projects in IA (intelligent artifice) are called Quest- Essence, Mind Pattern, Intellect, and Studio. His focus is on deterministic sequential models of digital emotion. 
During the 1960's, standing alone in the midst of a weed-'strewn field in this city, there was a weathered sign that read: -$25 reward for information leading to the arrest and conviction of anyone removing this sign." For whatever it's worth, the sign has long since disappeared. And so, for that matter, has the field. 
Incidentally, the Epimenides paradox should not be confused with the Nixonides paradox, first uttered by Nixonides the Cretin in A.D. 1974: "This statement is inoperative." Speaking of Epimenides, one of the most elegant variations on his paradox is the "Errata" section in a hypothetical book described by Beverly Rowe. It looks like this: 
\intro (vi) \endintro

Errata 
Page (vi): For Errata, read Erratum 

Closely related to the truly paradoxical sentences are those that belong to what I call the neurotic and healthy categories. A healthy sentence is one that, so to speak, practices what it preaches, whereas a neurotic sentence is one that says one thing while doing its opposite. Alan Auerbach has given us a good example in each category. His healthy sentence is: "Terse!" His neurotic sentence is: "Proper writing-and you've heard this a million times -avoids exaggeration." Here's a healthy one by Brad Shelton: "Fourscore and seven words ago, this sentence hadn't started yet." One of the jootsingest of sentences came from Carl Bender: 
The rest of this sentence is written in Thailand, on 
Consider a related sentence sent in by David Stork: "It goes without saying 



that ..." To which category does it belong? Perhaps it is a psychotic sentence. Pete Maclean contributed a puzzling one: "If the meanings of `true' and 'false' were switched, then this sentence wouldn't be false." I'm still scratching my head over what that means! Dan Krimm wrote to tell me: "I've heard that this sentence is a rumor." Linda Simonetti contributed the following example, "which actually is not a complete sentence, but merely a subordinate clause." Douglas Wolfe offered the following neurotic rule of thumb: "Never use the imperative, and it is also never proper to construct a sentence using mixed moods." David Moser reminded me of a slogan that the National Lampoon once used: "So funny it sells without a slogan!" Perry Weddle wrote, "I'm trying to teach my parrot to say, 'I don't understand a thing I say.' When I say it, it's viciously self-referential, but in his case?" Stephen Coombs pointed out that "A sentence may self-refer in the verb." My mother, Nancy Hofstadter, heard Secretary of State Alexander Haig describe a warning message to the Russians as "a calculated ambiguity that would be clearly understood". Yes, Sir! Jim Propp submitted a sequence of sentences that slide elegantly from the neurotically healthy to the healthily neurotic: 
(1) This sentence every third, but it still comprehensible. (2) This would easier understand fewer had omitted. (3) This impossible except context. (4) 4'33" attempt idea. \intro (5) \endintro

The penultimate sentence refers to John Cage's famous piece of piano music consisting of four minutes and 33 seconds of silence. The last sentence might well be an excerpt from The Wit and Wisdom of Spiro T. Agnew, although it is too short an excerpt to be sure. Propp also sent along the following healthy sentence, which was apparently inspired by his readings in the book Intelligence in Ape and Van, by David Premack: "By the 'productivity' of language, I mean the ability of language to introduce new words in terms of old ones." Philosopher Howard DeLong contributed what might be considered a neurotic syllogism: 
All invalid syllogisms break at least one rule. This syllogism breaks at least one rule. Therefore, this syllogism is invalid. 
Several readers pointed out phrases and jokes that have been making the rounds. D.A. Treissman, for instance, reminded me that "Nostalgia ain't what. it used to be." Henry Taves mentioned the delightful T-shirts adorned 



with statements such as "My folks went to Florida and all they brought back for me was this lousy T-shirt!" And John Fletcher described an episode of the television program Laugh-In a few years ago on which Joanne Worley sang, "I'm just a girl who can't say 'n . . .', 'n . . .', `n ...' ". John Healy wrote, "I used to think I was indecisive, but now I'm not so sure." I myself have a few contributions to this collection. A neurotic one is: "In this sentence, the concluding three words 'were left out'." Or is it neurotic? These things confuse me! In any case, a most healthy sentence is: "This sentence offers its reader(s) various alternatives/options that he or she (or they) is (are) free to accept and/or reject." And then there is the inevitable "This sentence is neurotic." The thing is, if it is neurotic, it practices what it preaches, so it's healthy and therefore cannot be neurotic-but then if it isn't neurotic, it's the opposite of what it claims to be, so it's got to be neurotic. No wonder it's neurotic, poor thing! Speaking of neurotic sentences, what about sentences with identity crises? These are, in some sense, the most interesting ones of all to me. A typical example is Dan Krimm's vaguely apprehensive question, "If I stated something else, would it still be me?" I thought this could be worded better, so I revised it slightly, as follows: "If I said something else, would it still be me saying it?" I still was not happy, so I wrote one more version: "In another world, could I have been a sentence about Humphrey Bogart?" When I paused to reflect on what I had done, I realized that in reworking Dan's sentence, I had tampered with its identity in the very way it feared. The question remained, however: Were all these variants really the same sentence, deep down? My last experiment along these lines was: "In another world, could this sentence have been Dan Krimm's sentence?" Clearly some readers were thinking along parallel lines, since John Atkins queried, "Can anyone explain why this would still be the same magazine without this query, and yet this would not be the same query without this word?" (Of course, just which word "this word" refers to is a little vague, but the idea is clear.) And Loul McIntosh, who works at a rehabilitation center for formerly schizophrenic patients, had a question connecting personal identity with self-referential sentences: "If I were you, who would be reading this sentence?" She then added: "That's what I get for working with schizophrenics." This brings me to Peter M. Brigham, M.D., who in his work ran across a severe case of literary schizophrenia: "You have, of course, just begun`reading the sentence that you have just finished reading." It's one of my favorites. Pursuing the slithery snake of self in his own way, Uilliam M. Bricken, Jr., wrote in: "If you think this sentence is confusing, then change one pig." Now, anyone can see that this doesn't make any sense at all. Surely what he meant was, "If you think this sentence is confusing, then roast one pig."don't ewe agree? By the by, if ewe think "Uilliam" is confusing, then roast one ewe. And while we're mentioning ewes, what's a nice word like "ewe" doing in a foxy paragraph like this? 



A while back, driving home late at night, I tuned in to a radio talk show about pets. A heated discussion was taking place about the relative merits of various species, and at one point the announcer mused, "If a dog had written this broadcast, he might have said that people are inferior because they don't wag their tails." This gave me paws for thought: What might this column have been like if it had been written by a dog? I can't say for sure, but I have a hunch it would have been about chasing squirrels. And it might have had a paragraph speculating about what this column would have been like if it had been written by a squirrel. 
\separator

I think my favorite of all the sent-in-ces was one contributed by Harold Cooper. He was inspired by Scott Kim's counterfactual self-referential question: "What would this sentence be like if π were 3?" His answer is shown in Figure 2-2. This, to me, exemplifies the meaning of the verb 



FIGURE 2-2. A counterfactual self-referential sentence, inspired by Harold Cooper and Scott Kim. 
"foots". The six-sided `o's represent the fact that the ratio of the circumference to the diameter of a hexagon is 3. Clearly, in Cooper's mind, if π were 3, why, what more natural conclusion than that circles would be hexagons! Who could ever think otherwise? I was intrigued by the fact that, as π´s value slipped to 3, not only did circles turn into hexagons, but also the interrogative mood slipped into the declarative mood. Remember that the question asked how the question itself would be in that strange subjunctive world. Would it lose its curiosity about itself and cease to be a question? I did not see why that personality trait of the sentence would be affected by the value of a. On the other hand, it seemed obvious to me that if π were 3, the antecedent of the conditional should no longer be subjunctive. In fact, rather than saying "if π were 3", it should say, "because π is 3" (or something to that effect). Putting my thoughts together, then, I came up with a slight variation on Cooper's sentence: "What is this sentence like, π being 3 (as usual)?" 



Several readers were interested in sentences that refer to the language they are in (or not in, as the case may be). An example is "If you spoke English, you'd be in your home language now." Jim Propp sent in a delightful pair of such sentences that need to be read together: 
Cette phrase se refere a elle-meme, mais d'une maniere peu evidente a la plupart des Americains. 
Plim glorkle pegram ut replat, trull gen ris clanter froat veb nup lamerack gla smurp Earthlings. 
If you do not understand the first sentence, just get a Martian friend to help you decode the second one. That will provide hints about the first. (I apologize for leaving off the proper Martian accent marks, but they were not available in this typeface.) 
\separator

Last January, I published several sentences by David Moser and mentioned that he had written an entire story consisting of self-referential sentences. Many readers were intrigued. I decided there could be no better way to conclude this column than to print David's story in its entirety. So here 'tis! 
This Is the Title of This Story, Which Is Also Found Several Times in the Story Itself This is the first sentence of this story. This is the second sentence. This is the title of this story, which is also found several times in the story itself. This sentence is questioning the intrinsic value of the first two sentences. This sentence is to inform you, in case you haven't already realized it, that this is a self-referential story, that is, a story containing sentences that refer to their own structure and function. This is a sentence that provides an ending to the first paragraph. This is the first sentence of a new paragraph in a self-referential story. This sentence is introducing you to the protagonist of the story, a young boy named Billy. This sentence is telling you that Billy is blond and blue-eyed and American and twelve years old and strangling his mother. This sentence comments on the awkward nature of the self-referential narrative form while recognizing the strange and playful detachment it affords the writer. As if illustrating the point made by the last sentence, this sentence reminds us, with no trace of facetiousness, that children are a precious gift from God and that the world is a better place when graced by the unique joys and delights they bring to it. This sentence describes Billy's mother's bulging eyes and protruding 



tongue and makes reference to the unpleasant choking and gagging noises she's making. This sentence makes the observation that these are uncertain and difficult times, and that relationships, even seemingly deep-rooted and permanent ones, do have a tendency to break down. Introduces. in this paragraph, the device of sentence fragments. A sentence fragment. Another. Good device. Will be used more later. This is actually the last sentence of the story but has been placed here by mistake. This is the title of this story, which is also found several times in the story itself. As Gregor Samsa awoke one morning from uneasy dreams he found himself in his bed transformed into a gigantic insect. This sentence informs you that the preceding sentence is from another story entirely (a much better one, it must be noted) and has no place at all in this particular narrative. Despite the claims of the preceding sentence, this sentence feels compelled to inform you that the story you are reading is in actuality "The Metamorphosis" by Franz Kafka, and that the sentence referred to by the preceding sentence is the only sentence which does indeed belong in this story. This sentence overrides the preceding sentence by informing the reader (poor, confused wretch) that this piece of literature is actually the Declaration of Independence, but that the author, in a show of extreme negligence (if not malicious sabotage), has so far failed to include even one single sentence from that stirring document, although he has condescended to use a small sentence fragment, namely, "When in the course of human events", embedded in quotation marks near the end of a sentence. Showing a keen awareness of the boredom and downright hostility of the average reader with regard to the pointless conceptual games indulged in by the preceding sentences, this sentence returns us at last to the scenario of the story by asking the question, "Why is Billy strangling his mother?" This sentence attempts to shed some light on the question posed by the preceding sentence but fails. This sentence, however, succeeds, in that it suggests a possible incestuous relationship between Billy and his mother and alludes to the concomitant Freudian complications any astute reader will immediately envision. Incest. The unspeakable taboo. The universal prohibition. Incest. And notice the sentence fragments? Good literary device. Will be used more later. This is the first sentence in a new paragraph. This is the last sentence in a new paragraph. This sentence can serve as either the beginning of the paragraph or the end, depending on its placement. This is the title of this story, which is also found several times in the story itself. This sentence raises a serious objection to the entire class of self- referential sentences that merely comment on their own function or placement within the story (e.g., the preceding four sentences), on the grounds that they are monotonously predictable, unforgivably self-indulgent, and merely serve to distract the reader from the real subject of this story, which at this point seems to concern strangulation and incest and who knows what other delightful 



topics. The purpose of this sentence is to point out that the preceding sentence, while not itself a member of the class of self-referential sentences it objects to, nevertheless also serves merely to distract the reader from the real subject of this story, which actually concerns Gregor Samsa's inexplicable transformation into a gigantic insect (despite the vociferous counterclaims of other well-meaning although misinformed sentences). This sentence can serve as either the beginning of a paragraph or the end, depending on its placement. This is the title of this story, which is also found several times in the story itself. This is almost the title of the story, which is found only once in the story itself. This sentence regretfully states that up to this point the self-referential mode of narrative has had a paralyzing effect on the actual progress of the story itself-that is, these sentences have been so concerned with analyzing themselves and their role in the story that they have failed by and large to perform their function as communicators of events and ideas that one hopes coalesce into a plot, character development, etc.-in short, the very raisons d'etre of any respectable, hardworking sentence in the midst of a piece of compelling prose fiction. This sentence in addition points out the obvious analogy between the plight of these agonizingly self-aware sentences and similarly afflicted human beings, and it points out the analogous paralyzing effects wrought by excessive and tortured self- examination. The purpose of this sentence (which can also serve as a paragraph) is to speculate that if the Declaration of Independence had been worded and structured as lackadaisically and incoherently as this story has been so far, there's no telling what kind of warped libertine society we'd be living in now or to what depths of decadence the inhabitants of this country might have sunk, even to the point of deranged and debased writers constructing irritatingly cumbersome and needlessly prolix sentences that sometimes possess the questionable if not downright undesirable quality of referring to themselves and they sometimes even become run-on sentences or exhibit other signs of inexcusably sloppy grammar like unneeded superfluous redundancies that almost certainly would have insidious effects on the lifestyle and morals of our impressionable youth, leading them to commit incest or even murder and maybe that's why Billy is strangling his mother, because of sentences just like this one, which have no discernible goals or perspicuous purpose and just end up anywhere, even in mid Bizarre. A sentence fragment. Another fragment. Twelve years old. This is a sentence that. Fragmented. And strangling his mother. Sorry, sorry. Bizarre. This. More fragments. This is it. Fragments. The title of this story, which. Blond. Sorry, sorry. Fragment after fragment. Harder. This is a sentence that. Fragments. Damn good device. The purpose of this sentence is threefold: (1) to apologize for the unfortunate and inexplicable lapse exhibited by the preceding paragraph; (2) to assure you, the reader, that it will not happen again; and (3) to 



reiterate the point that these are uncertain and difficult times and that aspects of language, even seemingly stable and deeply rooted ones such as syntax and meaning, do break down. This sentence adds nothing substantial to the sentiments of the preceding sentence but merely provides a concluding sentence to this paragraph, which otherwise might not have one. This sentence, in a sudden and courageous burst of altruism, tries to abandon the self-referential mode but fails. This sentence tries again, but the attempt is doomed from the start. This sentence, in a last-ditch attempt to infuse some iota of story line into this paralyzed prose piece, quickly alludes to Billy's frantic cover-up attempts, followed by a lyrical, touching, and beautifully written passage wherein Billy is reconciled with his father (thus resolving the subliminal Freudian conflicts obvious to any astute reader) and a final exciting police chase scene during which Billy is accidentally shot and killed by a panicky rookie policeman who is coincidentally named Billy. This sentence, although basically in complete sympathy with the laudable efforts of the preceding action-packed sentence, reminds the reader that such allusions to a story that doesn't, in fact, yet exist are no substitute for the real thing and therefore will not get the author (indolent goof-off that he is) off the proverbial hook. Paragraph. Paragraph. Paragraph. Paragraph. Paragraph. Paragraph. Paragraph. Paragraph. Paragraph. Paragraph. Paragraph. Paragraph. Paragraph. Paragraph. The purpose. Of this paragraph. Is to apologize. For its gratuitous use. Of. Sentence fragments. Sorry. The purpose of this sentence is to apologize for the pointless and silly adolescent games indulged in by the preceding two paragraphs, and to express regret on the part of us, the more mature sentences, that the entire tone of this story is such that it can't seem to communicate a simple, albeit sordid, scenario. This sentence wishes to apologize for all the needless apologies found in this story (this one included), which, although placed here ostensibly for the benefit of the more vexed readers, merely delay in a maddeningly recursive way the continuation of the by-now nearly forgotten story line. This sentence is bursting at the punctuation marks with news of the dire import of self-reference as applied to sentences, a practice that could prove to be a veritable Pandora's box of potential havoc, for if a sentence can refer or allude to itself, why not a lowly subordinate clause, perhaps this very clause? Or this sentence fragment? Or three words? Two words? One? Perhaps it is appropriate that this sentence gently and with no trace of condescension remind us that these are indeed difficult and uncertain times and that in general people just aren't nice enough to each other, and perhaps we, whether sentient human beings or sentient sentences, should just try harder. I mean, there is such a thing as free will, there has to be, and this sentence is proof of it! Neither this sentence nor you, the reader, is 



completely helpless in the face of all the pitiless forces at work in the universe. We should stand our ground, face facts, take Mother Nature by the throat and just try harder. By the throat. Harder. Harder, harder. Sorry. This is the title of this story, which is also found several times in the story itself. This is the last sentence of the story. This is the last sentence of the story. This is the last sentence of the story. This is. Sorry. 
Post Scriptum. As you can see, there is a vast amount of self-referential material out there in the world. To pick only the very best is a monumental task, and certainly a highly subjective one. I would like to include here some of the things that I had to omit from the second self-reference column with great regret, as well as some of the things that were sent in later, in response to it. First, though, I would like to mention an amusing incident. When Lee Sallows' self-documenting sentence was to be printed in the narrow columns of Scientific American, nobody remembered to tell the typesetters not to break any unhyphenated words. As luck would have it, two such breaks were introduced, yielding two spurious hyphens, thus spoiling (in a superficial sense) the accuracy of his construction. How subtly one can get snagged when self-reference is concerned! Paul Velleman sent me a copy of the front page of the Ithaca Journal, dated January 26, 1981, with a banner headline saying "Ex-hostages enjoy their privacy". He wrote, "I think it may be self-referent (and self-contradictory) in a different way than your other examples because the medium, positioning, and size of its printing are all necessary components of the contradiction." When I looked at the page, I simply saw nothing self-referential. I thought maybe I was supposed to look at the flip side, for some reason, but that had even less of interest. So I looked back at the headline, and suddenly it hit me: How can people "enjoy privacy" when it's being blared across the front page of newspapers across the nation? Along the same lines, soon thereafter I came across a photograph of Lady Di in tears, and in the caption her tears were explained this way: "Lady Di was apparently overcome by the strain of the impending royal wedding and having her every move in public watched by thousands. See story on page A20. Details on the royal honeymoon, page A7." John M. Lankford wrote me a long letter from Japan on self-reference, remarkably similar in some ways to the one from Flash gFiasco. The most memorable paragraph in his letter was the following one: 



Here in Japan, twice a week, I teach a little class in English for a group of university students-mainly graduate students in the sciences. I spent one class hour taking some of your sentences from the Scientific American article, writing them on the blackboard, and asking the students what they meant. The students had a fairly good command of written English, but they were poor in their command of idiom, quick verbal response, and, for want of a better term, "humor of the abstract". As I suspected, many of the sentences-perhaps the most interesting of them-die when ripped from their cultural context. I had quite a bit of difficulty getting across the idea that the pronoun "I" could refer to the sentence as well as to the writer of the sentence. Pronouns cause a lot of trouble in Japan. For example, when I ask someone, "Am I wearing a blue jacket?", they might frequently reply, "Yes, I am wearing a blue jacket." This confusion is easy in Japanese due to the relative lack of pronouns in ordinary speech. Of course you can imagine the extra layers of incomprehension that would arise in reading your sentences if the boundaries between "you" and "I" were rather vague. 
On a visit to Gettysburg, I read Abraham Lincoln's Gettysburg address, and for the first time its curious self-reference struck me: "The world will little note nor long remember what we say here." Lincoln had no way of knowing at the time, but this would turn out to be an extremely false sentence (if it is permissible to speak of degrees of falsity). In fact, that sentence itself is a very memorable one. While we're on presidential self-reference, listen to this self-descriptive remark by former President Ford: "I am the first to admit that I am no great orator or no person that got where I have gotten by any William Jennings Bryan technique." I guess that where Lincoln's sentence was extremely false, Ford's is extremely true. Here is a final self-referential sentence along presidential lines: 
If John F. Kennedy were reading this sentence, Lee Harvey Oswald would have missed. 
\separator

One of the best self-answering questions came up naturally in the course of a very brief telephone call I made to a restaurant one evening. It went this way: "May I help you?" to which I answered, "You've already helped me-by telling me that you're open today. Thank you. Bye!" And here's a "self-deferential" sentence by Don Byrd: "I am not as witty as my author." I received this anonymous letter in the mail: "I received this anonymous letter in the mail so I can't credit the author."-so I can't credit the author. I also received a request from someone living in Calgary, Alberta, whose name I forget (but if he's reading this, he'll know who he is) who wrote "This is my feeble way of attempting to get my name into print." I hope this satisfies him. And now a few miscellaneous examples by me, culled from a second wild binge of self-referential sentence-writing I engaged in not long ago. The first three involve translation issues. 


One me has translated at the foot of the letter of the French. 
Would not be anomalous if were in Italian. 
When one this sentence into the German to translate wanted, would one the fact exploit, that the word order and the punctuation already with the German conventions agree. 
How come this noun phrase doesn't denote the same thing as this noun phrase does? 
Every last word in this sentence is a grotesque misspelling of "towmatow". 
I don't care who wrote this sentence-whoever he is, he's a damn sexist! 
This analogy is like lifting yourself by your own bootstraps. 
Although this sentence begins with the word "because", it is false. 
Despite the fact that it opens like a two-pronged pitchfork-or rather, because of it-this sentence resembles a double-edged sword. 
This line from Shakespeare has delusions of grandeur. 
If writers were bakers, this sentence would be exactly a dozen words long. 
If this sentence had been on the previous page, this very moment would have occurred approximately 60 seconds ago. 
This sentence is helping to increase the likelihood of nuclear war by distracting you from the more serious concerns of the world and beguiling you with the trivial joys of self- reference. 
This sentence is helping to decrease the likelihood of nuclear war by chiding you for indulging in the trivial joys of self-reference and reminding you of the more serious concerns of the world. 
We mention "our gigantic nuclear arsenal" in order not to use it. 
The whole point of this sentence is to make clear what the whole point of this sentence is. 
This last one's bizarre circularity reminds me of the number P that I invented a couple of years ago. P is, for each individual, the number of 



minutes per month that that person spends thinking about the number P. For me, the value of P seems to average out at about 2. 1 certainly wouldn't want it to go much above that! I find it crosses my mind most often when I'm shaving. 
\separator

Dr. J. K. Aronson from Oxford, England, sent in some of the most marvelous discoveries. Here is one of his best: 
'T' is the first, fourth, eleventh, sixteenth, twenty-fourth, twenty-ninth, thirty- third.... 
The sentence never ends, of course. He also submitted a wonderful complementary pair that faked me out beautifully. His challenge to you is: Try deciphering the first before you read the second. 
I eee oai o ooa a e ooi eee o oe. 
Ths sntnc cntns n vwls nd th prcdng sntnc n cnsnnts. 
One that reminds me somewhat of Aronson's last sentence above is the following spoof on the ads that I believe you can still find in the New York subway, after all these years: 
f y cn rd ths, itn tyg h myxbl cd. 
By a remarkable coincidence, the remainder of Carl Bender's sentence "The rest of this sentence is written in Thailand, on" was discovered in, of all places, Bangkok, Thailand, by Gregory Bell, who lives there. He has luckily provided me with a perfect copy of it, so for all those who were dying of suspense, it is shown in Figure 2-3. One evening during a bad electrical storm, I got the following message on the computer from Marsha Meredith: 
I ]ion't be able to work at all tonight b]iecause of the w&atherBr/ I]i'm getting too many bad characters (as you can see). Ioo baw3d-I get spurious characters]i all over ]ithe place- talk totrrRBow,1F7U Marsha. 
FIGURE 2-3. The conclusion of Carl Bender's sentence fragment ("The rest of this sentence is written in Thailand, on"), discovered by Gregory Bell on a scrap of paper in Bangkok, Thailand , Translated it says: : "this sheet of paper and is in Thai". 



I wish she had had the patience to type more carefully, so that I could have understood what her problem was. The sentences having to do with identity in counterfactual worlds, such as Dan Krimm's and its alter egos, reminded me of a blurb by E. O. Wilson I read recently on Lewis Thomas' latest book: "If Montaigne had possessed a deep knowledge of twentieth- century biology, he would have been Lewis Thomas." Ah me, the flittering elf of self! And Banesh Hoffmann, in Relativity and Its Roots, has written: "How safe we would be from death by nuclear bomb had we been born in the time of Shakespeare." Sure, except we'd also all be long dead-unless, of course, the 24th-century doctors who will invent immortality pills had also been born in Shakespeare's time! The following self-referential poem just came to me one day: 
Twice five syllables, Plus seven, can't say much-but ... That's haiku for you. 
The genre of self-referential poetry-including haiku-was actually quite popular. Tom McDonald submitted this non-limerick: 
A very sad poet was Jenny Her limericks weren't worth a penny. In technique they were sound, Yet somehow she found Whenever she tried to write any, That she always wrote one line too many! 

Several people sent in complex poems of various sorts, and mentioned books of them, such as John Hollander's Rhyme's Reason, a collection of poems describing their own forms. 
\separator

Self-referential book titles are enjoying a mild vogue these days. Raymond Smullyan was one of the most enthusiastic explorers of the potential of this idea, using the titles What Is the Name of This Book? and This Book Needs No Title. Actually, I think Needs No Title would have said it more crisply, or maybe just No Title. Come to think of it, why not No, or even just plain ? (I hope you could tell that those blanks were in italics!) Other self-referential book titles I have collected include these: 



Forget all the rules you ever learned about graphic design. Including the ones in this book. Steal This Book Ban This Book Deduct This Book (How Not to Pay Taxes While Ronald Reagan Is President) Do You Think Mom Would Like This One? Dewey Decimal No. 510.46 FC H3 I Never Can Remember What It's Called The Great American Novel ISBN 0-943568-01-3 Self Referential Book Title The Top Book on the New York Times Bestseller List for the Past Ten Weeks Don't Go Overseas Until You've Read This Book Soon to Become a Major Motion Picture By Me, William Shakespeare (by Robert Payne) That Book with the Red Cover in Your Window Reviews of This Book 
Oh, by the way, some of these are fake, others are real. For example, the last one, Reviews of This Book, is just a fantasy of mine. I would love to see a book consisting of nothing but a collection of reviews of it that appeared (after its publication, of course) in major newspapers and magazines. It sounds paradoxical, but it could be arranged with a lot of planning and hard work. First, a group of major journals would all have to agree to run reviews of the book by the various contributors to the book. Then all the reviewers would begin writing. But they would have to mail off their various drafts to all the other reviewers very regularly so that all the reviews could evolve together, and thus eventually reach a stable state of a kind known in physics as a "Hartree-Fock self-consistent solution". Then the book could be published, after which its reviews would come out in their respective journals, as per arrangement. (A little more on this idea is given in the postscript to Chapter 16.) 
\separator

I chanced across two books devoted to the subject of indexing books. They are: A Theory of Indexing (by Gerald Salton) and Typescripts, Proofs, and Indexes (by Judith Butcher). Amazingly, neither one has an index. I also received a curious letter soliciting funds, which began this way: "Dear Friend: In these last months, I've been making a study of the money-raising letter as an art form ..." I didn't read any further. Aldo Spinelli, an Italian artist and writer, sent me some of his products. One, a short book called Loopings, has pages documenting their own word 



and letter counts in various complex ways, and includes at the end a short essay on various ways in which documents can tally themselves up or can mutually tally each other in twisty loops. Another, called Chisel Book, documents its own production, beginning with the idea, going through the finding of a publisher, making the layout, designing the cover, printing it, and so on. Ashleigh Brilliant is the inventor of a vast number of aphorisms he calls "potshots", many of which have become very popular phrases in this country. For some reason, he has a self-imposed limit of seventeen words per potshot. A few typical potshots (all taken from his four books listed in the Bibliography) are: 
What would life be, without me? 
As long as I have you, I can endure all the troubles you inevitably bring. 
Remember me? I'm the one who never made any impression on you. 
Why does trouble always come at the wrong time? 
Due to circumstances beyond my control, I am master of my fate and captain of my soul. 
Although strictly speaking these are not self-referential sentences, they are all admirable examples of how the world constantly tangles with itself in multifarious self-undermining ways, and as such, they definitely belong in this chapter. As a matter of fact, I would like to take this occasion to announce that Ashleigh Brilliant is the 1984 recipient of the last annual Nobaloney Prize for Aphoristic Eloquence. The traditional Nobaloney ceremony, involving the awarding of a $1,000,000 cash prize two minutes before the recipient's decapitation, has been waived, at Mr. Brilliant's request. There are other books containing much of interest to the self-reference addict. I would particularly recommend the recent More on Oxymoron, by Patrick Hughes, as well as the earlier Vicious Circles and Infinity, by Hughes and George Brecht. Also in this category are three thin volumes on Murphy's Law, compiled by Arthur Bloch. Murphy's Law, of course, is the one that says, "If anything can go wrong, it will", although when I first heard of it, it was called the "Fourth Law of Thermodynamics". O'Toole's Commentary on Murphy's Law is: "Murphy was an optimist." Goldberg's Commentary thereupon is: "O'Toole was an optimist." And finally, there is Schnatterly's Summing Up: "If anything can't go wrong, it will." My own law, "Hofstadter's Law", states: "It always takes longer than you think it will take, even if you take into account Hofstadter's Law." Despite being its enunciator, I never seem to be able to take it fully into account in 



budgeting my own time. To help me out, therefore, my friend Don Byrd came up with his own law that I have taken to heart: 
Byrd's Law: 
It always takes longer than you think it will take, even if you take into account Hofstadter's Law. 
Unfortunately, Byrd himself seems unable to take this law into account. 



On Viral Sentences and Self-Replicating Structures January, 1983 
TWO years ago, when I first wrote about self-referential sentences, I was hit by an avalanche of mail from readers intrigued by the phenomenon of self-reference in its many different guises. I had the chance to print some of those responses one year ago, and that column then triggered a second wave of replies. Many of them have cast self-reference in new light of various sorts. In this column, I would like to describe the ideas of several people, two of whom responded to my initial column with remarkably similar letters: Stephen Walton of New York City and Donald R. Going of Oxon Hill, Maryland. Walton and Going saw self-replicating sentences as similar to virusessmall objects that enslave larger and more self-sufficient "host" objects, getting the hosts by hook or by crook to carry out a complex sequence of replicating operations that bring new copies into being, which are then free to go off and enslave further hosts, and so on. "Viral sentences", as Walton called them, are "those that seek to obtain their own reproduction by commandeering the facilities of more complex entities". Both Walton and Going were struck by the perniciousness of such sentences: the selfish way in which they invade a space of ideas and, merely by making copies of themselves all over the place, manage to take over a large portion of that space. Why do they not manage to overrun all of that idea-space? A good question. The answer should be obvious to students of evolution: competition from other self-replicators. One type of replicator seizes a region of the space and becomes good at fending off rivals; thus a "niche" in idea-space is carved out. This idea of an evolutionary struggle for survival by self-replicating ideas is not original with Walton or Going, although both had fresh things to say on it. The first reference I know of to this notion is in a passage by neurophysiologist Roger Sperry in an article he wrote in 1965 called "Mind, Brain, and Humanist Values". He says: "Ideas cause ideas and help evolve 

\bye
